\documentclass[11pt]{article}
\usepackage[margin=24mm]{geometry}
\usepackage[T1]{fontenc}
\usepackage{amsmath,amssymb,amsthm,mathtools,mathrsfs}
\usepackage[hidelinks]{hyperref}
\allowdisplaybreaks
\setlength{\emergencystretch}{3em}
\title{Original theta closure and arithmetic observation: complete literature receivers}
\author{}
\date{20 September 2026}
\begin{document}\maketitle
This edition proves three exact uses of the literature in the original
Split-Zero cohomology programme: the original Gamma evaluation kernel and
its three-vector arithmetic observation constraints; closure of the full
theta source in its specified strong topology; and reconstruction of the
actual arithmetic Gram from a single coefficient sequence, retaining the
finite Jacobi boundary correction. Each part contains its full derivation
and identifies the human source at the point of use. No sign is assigned
to an unevaluated projected current. These results do not close the RH
programme or assert existence of an off-line zero.

% BEGIN INDEXED LITERATURE EXACT RECEIVERS 20260920

\clearpage\begingroup
\part*{Exact Gamma kernel and the actual joint observed Gram}
\newtheorem{lgammatheorem}{Theorem}
\newtheorem{lgammalemma}[lgammatheorem]{Lemma}
\newtheorem{lgammaproposition}[lgammatheorem]{Proposition}
\newcommand{\lgammaC}{\mathbb C}
\newcommand{\lgammaR}{\mathbb R}
\newcommand{\lgammanorm}[1]{\left\lVert#1\right\rVert}
\newcommand{\lgammaim}{\operatorname{im}}
\newcommand{\lgammaremm}{\operatorname{rem}}



This is a standalone mathematical derivation.  All source masses,
coordinates, physical units, observation maps, original arithmetic boundary
vectors and affine relation spaces remain present.  The new finite Gamma
kernel formula is transported to the arithmetic source using the already
proved whole-polynomial comparison in CAI16--18 and OB13,18
\cite{lgamma:CAI,lgamma:OB}.  The receiving object is the actual three-vector Gram from
HG6--7 and HR2--5 \cite{lgamma:HG,lgamma:HR}, not a newly selected Gamma boundary triple.
The resulting constraints retain their positive-semidefinite remainders
and every complex cross term.  No projected current sign is assigned by
discarding an unevaluated remainder.

\section{The original Gamma measure and its exact norms}

Put
\begin{equation}\label{lgamma:eq:sigma}
 d\sigma(y)=\frac{|\Gamma(\tfrac14+iy/2)|^2}{2\pi}\,dy,
 \qquad M_\Gamma=\sqrt{2\pi}.
\end{equation}
The following calculation both fixes the mass and proves the polynomial
identities used below; no change of measure is made.
Euler's beta integral, with $a>0$, gives
\[
 |\Gamma(a+it)|^2
 =\Gamma(2a)\int_{\lgammaR}e^{itu}(2\cosh(u/2))^{-2a}\,du.
\]
Indeed this is the substitution $x=e^u$ in the beta integral on
$(0,\infty)$.  The function of $u$ is integrable.  The gamma product is
integrable as well, since Stirling's formula gives exponential decay
$e^{-\pi|t|}$ times a power.  Fourier inversion and evenness give
\[
 \int_{\lgammaR}|\Gamma(a+it)|^2e^{i\xi t}\,dt
 =2\pi\Gamma(2a)(2\cosh(\xi/2))^{-2a}.
\]
At $a=1/4$, substituting $y=2t$ in this integral, with the Jacobian
retained, proves
\begin{equation}\label{lgamma:eq:transform}
 \int_{\lgammaR}e^{i\xi y}\,d\sigma(y)
 =M_\Gamma(\cosh\xi)^{-1/2}.
\end{equation}
This substitution proves an integral identity and does not redefine the
measure or the polynomial coordinate.  In particular
$\sigma(\lgammaR)=M_\Gamma$.  The same gamma decay proves exponential
integrability for real $s$ with $|s|<\pi/2$.  Analytic continuation in
$|\Im\xi|<\pi/2$, with the branch positive at zero, therefore gives
\begin{equation}\label{lgamma:eq:laplace}
 \int_{\lgammaR}e^{sy}\,d\sigma(y)=M_\Gamma(\cos s)^{-1/2}.
\end{equation}

Define the monic polynomials in the original variable by
\begin{equation}\label{lgamma:eq:gen}
 F(y,t)=\sum_{n\ge0}p_n^\sigma(y)\frac{t^n}{n!}
 =(1+t^2)^{-1/4}e^{y\arctan t}.
\end{equation}
For sufficiently small real $t,u$, the cosine addition identity and
\eqref{lgamma:eq:laplace} give
\[
 \int_{\lgammaR}F(y,t)F(y,u)\,d\sigma(y)=M_\Gamma(1-tu)^{-1/2}.
\]
Exponential integrability justifies every derivative at $t=u=0$.
Comparing coefficients proves
\begin{equation}\label{lgamma:eq:norms}
 \int p_n^\sigma p_m^\sigma\,d\sigma
 =\delta_{nm}\gamma_n,
 \qquad \gamma_n=M_\Gamma n!(\tfrac12)_n.
\end{equation}
Also $(1+t^2)\partial_tF=(y-t/2)F$, so
\begin{equation}\label{lgamma:eq:monicrec}
 p_0^\sigma=1,\quad p_1^\sigma=y,\qquad
 p_{n+1}^\sigma=yp_n^\sigma-n(n-\tfrac12)p_{n-1}^\sigma.
\end{equation}
Thus the Gamma polynomial formulas are proved here directly.  Their
existing programme occurrence is CAI19 \cite{lgamma:CAI}; no result from CCM
is used in this derivation.

\section{Exact finite evaluation of the imaginary-point kernel}

\begin{lgammatheorem}\label{lgamma:thm:kernel}
For $N\ge0$, put $m_N=\lceil N/2\rceil$ and
\[
 \tau_m=\frac{(\tfrac12)_m(\tfrac34)_m}{m!(\tfrac14)_m}.
\]
For precisely the measure \eqref{lgamma:eq:sigma},
\begin{equation}\label{lgamma:eq:exactkernel}
 \boxed{K_N(i,i):=\sum_{n=0}^N\frac{|p_n^\sigma(i)|^2}{\gamma_n}
       =\frac{N+1}{M_\Gamma}\tau_{m_N}.}
\end{equation}
In particular, the exact reciprocal scalar is
\begin{equation}\label{lgamma:eq:beta}
 \boxed{\beta_N=K_N(i,i)^{-1}
 =\frac{\sqrt{2\pi}}{(N+1)\tau_{\lceil N/2\rceil}}.}
\end{equation}
The following explicit finite-degree bounds also hold:
\begin{align}
 K_N(i,i)&\le\frac{2N+1}{M_\Gamma}&& (N\ge0),\label{lgamma:eq:linear}\\
 K_N(i,i)&\le\frac{18N+9}{10M_\Gamma}&& (N\ge2).\label{lgamma:eq:betterlinear}
\end{align}
\end{lgammatheorem}
\begin{proof}
Write $p_n^\sigma(i)=i^na_n$.  Equation \eqref{lgamma:eq:monicrec} proves
\[
 a_0=a_1=1,\qquad
 a_{n+1}=a_n+n(n-\tfrac12)a_{n-1}>0.
\]
Inductively,
\begin{equation}\label{lgamma:eq:alternating}
 \frac{a_n}{a_{n-1}}=
 \begin{cases}n,&n\text{ odd},\\n-\tfrac12,&n\text{ even}.
 \end{cases}
\end{equation}
For the induction step, $a_n=na_{n-1}$ at odd $n$ gives
$a_{n+1}=(n+1/2)a_n$, while $a_n=(n-1/2)a_{n-1}$ at even $n$
gives $a_{n+1}=(n+1)a_n$.  The initial ratio is one.  Multiplication gives
\[
 a_{2m}=4^m(\tfrac12)_m(\tfrac34)_m,
 \qquad a_{2m+1}=(2m+1)a_{2m}.
\]
Let $u_n=a_n/\sqrt{\gamma_n}$.  Grouping the even and odd factors
in $(2m)!$ and $(1/2)_{2m}$ gives
\begin{equation}\label{lgamma:eq:terms}
 M_\Gamma u_{2m}^2=\tau_m,\qquad
 M_\Gamma u_{2m+1}^2=\frac{2m+1}{2m+\tfrac12}\tau_m,
 \qquad
 \frac{\tau_{m+1}}{\tau_m}
 =\frac{(m+\tfrac12)(m+\tfrac34)}{(m+1)(m+\tfrac14)}.
\end{equation}
The partial sums satisfy exactly
\[
 M_\Gamma K_{2m}=(2m+1)\tau_m,
 \qquad M_\Gamma K_{2m+1}=(2m+2)\tau_{m+1}.
\]
To prove this, start at $K_0=M_\Gamma^{-1}$.  Adding the odd term
in \eqref{lgamma:eq:terms} to the asserted even sum gives
$(2m+1)(2m+3/2)\tau_m/(2m+1/2)=(2m+2)\tau_{m+1}$.
Adding the next even term gives $(2m+3)\tau_{m+1}$.  This proves
\eqref{lgamma:eq:exactkernel} at every degree.

For the linear estimates, the exact original-coordinate coefficients are
\[
 b_n=\sqrt{\gamma_{n+1}/\gamma_n}=\sqrt{(n+1)(n+\tfrac12)},
 \qquad u_{n+1}=\frac{u_n+b_{n-1}u_{n-1}}{b_n}.
\]
For $n\ge1$,
\[
 b_n^2-(1+b_{n-1})^2
 =2n-\tfrac12-2\sqrt{n(n-\tfrac12)}>0,
\]
since both compared quantities are positive and
$(2n-1/2)^2-4n(n-1/2)=1/4$.  Therefore
\begin{equation}\label{lgamma:eq:pairmax}
 u_{n+1}\le\frac{1+b_{n-1}}{b_n}\max(u_{n-1},u_n)
 <\max(u_{n-1},u_n).
\end{equation}
The initial squares are $1/M_\Gamma$ and $2/M_\Gamma$, proving
\eqref{lgamma:eq:linear}.  The next squares are $3/(2M_\Gamma)$ and
$9/(5M_\Gamma)$.  Starting \eqref{lgamma:eq:pairmax} from this pair bounds
every term from index three on by $9/(5M_\Gamma)$.  Summing gives
$M_\Gamma K_N\le1+2+3/2+(N-2)9/5$, proving
\eqref{lgamma:eq:betterlinear}.
\end{proof}

\section{The disk estimate and the exact coefficient-to-observation map}

The imaginary-point domination is the argument in Christian Berg and
Ryszard Szwarc's Section 3 \cite{lgamma:BS}.  Their source starts with mass one
and $b_0=1$; the following proof retains the original finite mass and
recurrence coefficients.  Their label \texttt{thm:lambda} is a lemma,
not a theorem.

\begin{lgammalemma}\label{lgamma:lem:disk}
Let $m$ be a nonzero finite symmetric positive measure on $\lgammaR$, with
moments through order $2N$ and positive definite original monomial Gram
$H_y=(\int y^{j+\ell}\,dm)_{j,\ell=0}^N$.  If $P_n$ are its
orthonormal polynomials with positive leading coefficient and
$K_N^m(i,i)=\sum_{n=0}^N|P_n(i)|^2$, then
\begin{equation}\label{lgamma:eq:diskgram}
 |P_n(z)|\le|P_n(i)|\quad(|z|\le1),\qquad
 H_y\succeq K_N^m(i,i)^{-1}I.
\end{equation}
\end{lgammalemma}
\begin{proof}
Symmetry gives parity $P_n(-y)=(-1)^nP_n(y)$.  Expanding $yP_n$ in
the orthonormal basis shows that its coefficients below index $n-1$
vanish: pair instead $P_n$ with $yP_j$ of degree at most $n-1$.
The coefficient at index $n$ vanishes by parity.  The leading-coefficient
ratio gives $b_n>0$ at index $n+1$; symmetry of multiplication gives
$b_{n-1}$ at index $n-1$.  Consequently
\[
 yP_n=b_nP_{n+1}+b_{n-1}P_{n-1}.
\]
Starting at $P_0=m(\lgammaR)^{-1/2}$ and
$P_1(y)=m(\lgammaR)^{-1/2}y/b_0$, induction gives
$P_n(i)=i^nv_n$ with $v_n>0$ and
$v_{n+1}=(v_n+b_{n-1}v_{n-1})/b_n$.
The triangle inequality in the same recurrence proves disk domination
by induction.  Degree zero is immediate.

For $f(y)=\sum_{j=0}^Nc_jy^j=\sum_{n=0}^Nd_nP_n(y)$,
Cauchy--Schwarz and disk domination give on $|z|=1$
\[
 |f(z)|^2\le K_N^m(i,i)\sum|d_n|^2
 =K_N^m(i,i)c^*H_yc.
\]
Integrating over the circle gives exactly
$\sum|c_j|^2\le K_N^m(i,i)c^*H_yc$ by Fourier orthogonality.
No source mass or coordinate was changed.
\end{proof}

Let the original complex coordinate be $S=c+iy$, with the actual real
$c$ retained.  For $F(S)=\sum a_jS^j$ and $f(y)=F(c+iy)=\sum b_ry^r$,
binomial expansion gives $b=T_N(c)a$, with
\begin{equation}\label{lgamma:eq:T}
 T_{rj}=\binom jr c^{j-r}i^r,\qquad
 (T^{-1})_{rj}=\binom jr(-c)^{j-r}i^{-j}\quad(r\le j),
\end{equation}
and both entries zero if $r>j$.  The inverse formula follows by
expanding $((S-c)/i)^j$.  In particular,
\begin{equation}\label{lgamma:eq:Tgram}
 \det T=i^{N(N+1)/2},\quad H_S=T^*H_yT,\quad
 H_S^{-1}=T^{-1}H_y^{-1}T^{-*},\quad \det H_S=\det H_y.
\end{equation}
Here $H_S$ is the literal Gram of $1,S,\ldots,S^N$, and
$T^{-*}=(T^{-1})^*$.  Thus $H_y\succeq\beta I$ implies
$H_S\succeq\beta T^*T$, not an unsupported bound by $\beta I$.

Let $J_{S,N}$ be the actual native coefficient-to-class matrix and let
$\Lambda$ be the original observation onto its image in its fixed frame.
All physical-unit and divided-jet factors occurring in $J_{S,N}$ remain
there.  Put
\begin{equation}\label{lgamma:eq:A}
 A_{S,N}=\Lambda J_{S,N},\quad
 A_{y,N}=A_{S,N}T^{-1},\quad
 \mathcal B_N=A_{S,N}T^{-1}T^{-*}A_{S,N}^*.
\end{equation}
Whenever the source cutoff maps onto the class space, $A_{S,N}$ has full
row rank on the observation image, so $\mathcal B_N>0$.

\begin{lgammalemma}\label{lgamma:lem:min}
For a positive definite coefficient Gram $H$ and full-row-rank $A$,
\begin{equation}\label{lgamma:eq:min}
 \min_{Aa=w}a^*Ha=w^*(AH^{-1}A^*)^{-1}w,\qquad
 a_{\min}=H^{-1}A^*(AH^{-1}A^*)^{-1}w.
\end{equation}
For the unchanged Gamma source at $S=c+iy$ this gives
\begin{equation}\label{lgamma:eq:Qsigma}
 \boxed{Q_N^\sigma=(A_{S,N}(H_S^\sigma)^{-1}A_{S,N}^*)^{-1}
 \succeq\beta_N\mathcal B_N^{-1}.}
\end{equation}
\end{lgammalemma}
\begin{proof}
The proposed minimizer has image $w$.  Every other lift is $a_{\min}+h$
with $Ah=0$, and $Ha_{\min}=A^*(AH^{-1}A^*)^{-1}w$ is orthogonal
to $h$.  Its energy is the minimum in \eqref{lgamma:eq:min} plus $h^*Hh$.
Lemma \ref{lgamma:lem:disk} and Theorem \ref{lgamma:thm:kernel} give
$(H_y^\sigma)^{-1}\preceq\beta_N^{-1}I$.  Conjugate by
$A_{S,N}T^{-1}$ and invert to obtain \eqref{lgamma:eq:Qsigma}.  Inversion
reverses order because conjugating $0<X\preceq Y$ by $Y^{-1/2}$
gives a positive matrix with eigenvalues at most one, whose inverse has
eigenvalues at least one.
\end{proof}

For clarity, if $V_S,V_y$ are the literal derivative-jet matrices at
$s_\ell=c+iy_\ell$ with row orders $d$, the chain rule gives
\begin{equation}\label{lgamma:eq:jets}
 V_yT=D V_S,\qquad D_{(\ell,d),(\ell,d)}=i^d.
\end{equation}
It also holds for corresponding divided derivatives.  Thus if the
actual map is $J_{S,N}=U_{\rm phys}V_S$, its transform is
$J_{S,N}T^{-1}=U_{\rm phys}D^{-1}V_y$: the physical unit is not removed,
and $D^{-1}$ is not moved through it.  For square jet matrices,
\[
 \det V_S=i^{N(N+1)/2-\sum d}\det V_y.
\]
This keeps the determinant phase, including the case of zero determinants.

\section{The native simple-quartet domain, weight and exact window}

We now use exactly the simple-quartet domain of CAI and OB
\cite{lgamma:CAI,lgamma:OB}: a stipulated full simple quartet of zeros of $g=2\xi$,
\begin{equation}\label{lgamma:eq:nativedomain}
 \rho=\tfrac12+\delta+i\gamma,\quad 0<\delta<\tfrac12,\quad\gamma>2,
 \quad \mathcal R=\{\rho,\bar\rho,1-\rho,1-\bar\rho\},
 \quad h(s)=\prod_{\zeta\in\mathcal R}(s-\zeta).
\end{equation}
Here $\xi(s)=\tfrac12s(s-1)\pi^{-s/2}\Gamma(s/2)\zeta(s)$.
This domain does not assert existence of an off-line quartet.
All statements below concern the original objects on that stipulated
domain, without adding a zero, period or phase assumption.
Retain
\begin{equation}\label{lgamma:eq:nativeweight}
 v_h=g/h,\quad w_h(y)=\frac{|v_h(1/2+iy)|^2}{2\pi},\quad
 dm_k(y)=w_h^{*k}(y)\,dy,\quad
 \mu_h=\int w_h(y)\,dy.
\end{equation}
The arithmetic mass is $\mu_h^k$, not $M_\Gamma$ and not
$M_\Gamma^k$.  Fix the original period and its admitted observation map.
The original parameters and polynomial are
\begin{equation}\label{lgamma:eq:chik}
 \begin{gathered}
 k=4\ell+1\ge9,\quad c=k/2,\quad q=(k+1)^2,\quad S=c+iy,\\
 \chi_k(S)=\prod_{a,b=0}^k
 [S-c-(2a-k)\delta-i(2b-k)\gamma],\qquad E_k=\lgammaC[S]/(\chi_k).
 \end{gathered}
\end{equation}
The roots are distinct since $\delta,\gamma>0$.  The exact remainder map
is onto for $N\ge q-1$ and has kernel $\chi_k\lgammaC[S]_{\le N-q}$.
The complete original cutoff window is
\begin{equation}\label{lgamma:eq:window}
 q-1\le N\le2q,
 \qquad N_*=2q+2.
\end{equation}
The four return endpoints are $q-1,q,2q-1,2q$; the adjacent-current sum
uses $N=q-1,\ldots,2q-1$ and weights $1,2,\ldots,2,1$ of total $2q$.
The larger $N_*$ is exactly the CAI comparison degree used in OB13 and
FW7; it is not a replacement window.

\paragraph{The unit and the exact physical presentation.}
The original local Taylor unit is $\upsilon_h=j_hv_h$, invertible
because the full simple zeros have been divided out.  In the physical
tensor presentation the original map is
\[
 \eta:E_k\longrightarrow E_h^{\otimes k},\qquad
 \eta[P]=\upsilon_h^{\otimes k}[P(s_1+\cdots+s_k)],
 \qquad E_h=\lgammaC[s]/(h).
\]
It is injective: all grid centres in \eqref{lgamma:eq:chik} occur as sums of
the retained quartet.  For given $a,b$, choose $x=\max(0,a+b-k)$
and quartet multiplicities $x,a-x,b-x,k-a-b+x$; these are nonnegative
and have the required two marginal counts.  Vanishing at all tensor
points, after multiplying
by the invertible unit, is equivalent to divisibility by $\chi_k$.
Thus the actual physical coefficient-to-class map is
$J_{S,N}=\eta\circ\lgammaremm_{\chi_k}$, with codomain the original image
$\eta(E_k)$, whenever this is the native frame used.  Its observation
is the retained period map, whose included form is
$P_k\Pi^{\otimes k}$ on that image.  In CAI's literal remainder
presentation the identical composite is written
\[
 \jmath\Lambda_{\rm rem}
 =P_k\Pi^{\otimes k}U^{\otimes k}\iota,
 \qquad U=M_{j_hv_h},\qquad \iota[P]=[P(\textstyle\sum s_i)].
\]
These are equal maps, not a deletion of $U$.  If $R$ is the fixed
matrix of $\eta$ between the two retained class frames, their exact
relations are
\begin{equation}\label{lgamma:eq:unitframe}
 J_{S,N}=R J_{{\rm rem},N},\quad
 \Lambda=\Lambda_{\rm rem}R^{-1},\quad
 G_N=R^{-*}G_{{\rm rem},N}R^{-1},\quad
 a_i=R a_{i,\rm rem},\quad v=Rv_{\rm rem}.
\end{equation}
Consequently $A_{S,N}$ and every observed vector below are unchanged.
This retains the original unit, including its actual phase.  In
particular the unit matrix $\operatorname{diag}(a,\bar a,\bar a,a)$
of NC1--3 \cite{lgamma:NC} is not set to the identity.  Neither NC's
proper-source kernel nor its independent target minimum is substituted
for the present observation fibre.

\paragraph{The established CAI constants in the simple domain.}
The original constants from CAI16--18, with common multiplicity one,
are exactly
\begin{equation}\label{lgamma:eq:CAIconstants}
 \begin{gathered}
 \alpha=\pi/2,\quad p=7/2,\quad B_h=50,\quad M=25,\\
 \vartheta_h=\int_{-1}^1w_h(y)\,dy>0,\quad
 M_\alpha=\int_\lgammaR e^{\alpha y}w_h(y)\,dy\in(0,\infty),\\
 c_\Gamma=\sqrt2e^{-7/6},\qquad C_\Gamma=\sqrt{2\sqrt5}\,e^{2/3},\\
 a_k=\frac{c_h\vartheta_h^{k-3}e^{-\alpha(k-3)}(k-2)^{-50}}
              {C_\Gamma2^{25}},\qquad
 A_k=\frac{C_h^{\rm tilt}}{c_\Gamma}k^{9/2}M_\alpha^{k-1},\\
 D_n=[1+4(n+25)^2]^{25},\qquad
 \ell_k=\frac{a_k}{[1+4(2q+2+25)^2]^{25}},\qquad u_k=A_k.
 \end{gathered}
\end{equation}
Here $c_h$ is the original TW7 lower-envelope constant and
$C_h^{\rm tilt}$ the original BT12 compact/tail maximum as retained in
CAI16--17, not newly selected constants.  The established density
inequalities in that provider are, for every convolution order $b\ge3$,
\begin{align}
 w_h^{*b}(y)&\ge c_h\vartheta_h^{b-3}
 e^{-\alpha(|y|+b-3)}(b-2+|y|)^{-50},\label{lgamma:eq:CAIlowdens}\\
 w_h^{*b}(y)&\le bC_h^{\rm tilt}M_\alpha^{b-1}
 e^{-\alpha|y|}(1+|y|/b)^{-7/2},\label{lgamma:eq:CAIupdens}\\
 c_\Gamma e^{-\alpha|y|}(1+|y|)^{-1/2}
 &\le\frac{d\sigma}{dy}
 \le C_\Gamma e^{-\alpha|y|}(1+|y|)^{-1/2}.
 \label{lgamma:eq:CAIgammadens}
\end{align}
These quoted full-density results are the native analytic input.  For
completeness their passage to the exact polynomial receiver is as follows.
The elementary inequalities
$b-2+|y|\le(b-2)(1+|y|)$,
$(1+|y|)^{50}\le2^{25}(1+y^2)^{25}$, and
$1+|y|/b\ge(1+|y|)/b$ give
\[
 a_b(1+y^2)^{-25}\,d\sigma\le w_h^{*b}(y)\,dy\le A_b\,d\sigma.
\]
The original Gamma recurrence gives
$\lgammanorm{yP}_\sigma\le2(n+1)\lgammanorm P_\sigma$ at degree $n$:
its raising and lowering parts have respective norms at most $n+1$
and $n$ on that coefficient span.  Iterating at successive degrees gives
$\lgammanorm{y^jP}_\sigma\le[2(n+25)]^j\lgammanorm P_\sigma$ for $0\le j\le25$.
Expand $(1+y^2)^{25}$ to obtain
\[
 \int|P|^2(1+y^2)^{25}\,d\sigma\le D_n\lgammanorm P_\sigma^2.
\]
Cauchy--Schwarz with the reciprocal weights gives
$\int|P|^2(1+y^2)^{-25}\,d\sigma\ge D_n^{-1}\lgammanorm P_\sigma^2$.
Thus, for every complex polynomial of degree at most $N_*=2q+2$,
\begin{equation}\label{lgamma:eq:nativecomparison}
 \boxed{\ell_k\lgammanorm{P}_{\sigma,c}^{2}
       \le\lgammanorm{P}_{m_k,c}^{2}
       \le u_k\lgammanorm{P}_{\sigma,c}^{2}.}
\end{equation}
Here the subscript $c$ means that the same original polynomial is
evaluated at $S=c+iy$.  Its degree is unchanged by the explicitly
proved map $T_N(c)$.  This is CAI18 at $N_*$, exactly as used in
OB13 and OB18.  It does not identify $w_h^{*k}$ with $\sigma$ or with
$\sigma^{*k}$.

\section{Identical-fibre transfer to the native observed form}

Let $H_{S,N}^{\rm ar}$ be the full original coefficient Gram for
\eqref{lgamma:eq:nativeweight}, and let $G_N^{\rm ar}$ be the attained class
form at degree $N$, including every original relation.  In the retained
class frame,
\[
 (G_N^{\rm ar})^{-1}
 =J_{S,N}(H_{S,N}^{\rm ar})^{-1}J_{S,N}^*,\quad
 Q_N^{\rm ar}=(\Lambda(G_N^{\rm ar})^{-1}\Lambda^*)^{-1}.
\]
Composition therefore gives exactly
$Q_N^{\rm ar}=(A_{S,N}(H_{S,N}^{\rm ar})^{-1}A_{S,N}^*)^{-1}$.
All these matrices exist for \eqref{lgamma:eq:window}; the densities have
positive norm on nonzero polynomials, the remainder map is onto, and
the observation is taken onto its image.

\begin{lgammatheorem}[Native observed lower form]\label{lgamma:thm:nativelower}
At every original cutoff $q-1\le N\le2q$,
\begin{equation}\label{lgamma:eq:samefibre}
 \ell_kQ_N^\sigma\preceq Q_N^{\rm ar}\preceq u_kQ_N^\sigma,
\end{equation}
where the Gamma quotient uses the identical $A_{S,N}$, period,
physical-unit factors and observation frame.  Consequently
\begin{equation}\label{lgamma:eq:nativelower}
 \boxed{Q_N^{\rm ar}\succeq
 \ell_k\beta_N\left[
 \Lambda J_{S,N}T_N(k/2)^{-1}T_N(k/2)^{-*}J_{S,N}^*\Lambda^*
 \right]^{-1}.}
\end{equation}
The scalar $\beta_N$ is the exact formula \eqref{lgamma:eq:beta}.
\end{lgammatheorem}
\begin{proof}
Fix an arbitrary observation vector $w$.  The candidate set in both
minima is exactly
\[
 \{P\in\lgammaC[S]_{\le N}:A_{S,N}\operatorname{coeff}_S(P)=w\}.
\]
It is nonempty and is the same affine fibre, with every relation and
kernel lift included.  Apply \eqref{lgamma:eq:nativecomparison} to every
candidate.  Taking infima proves the lower inequality.  For the upper
inequality evaluate the arithmetic norm at the Gamma minimizer;
Lemma \ref{lgamma:lem:min} proves its existence and its exact norm.  This
proves \eqref{lgamma:eq:samefibre} for every $w$, hence in Loewner order.
Combine its lower inequality with \eqref{lgamma:eq:Qsigma} and retain
$\mathcal B_N$ from \eqref{lgamma:eq:A}.  This proves \eqref{lgamma:eq:nativelower}.
It is precisely the full-matrix version of OB18's same-class minimum
argument.  No extra large-period or conductor-corner hypothesis from
OB14 is needed or imposed.
\end{proof}

If the observation image has dimension zero, all observed vectors and
forms below are zero and the same assertions hold with empty-matrix
conventions.  Otherwise the inverses above are ordinary positive
definite inverses.

\section{The actual arithmetic boundary triple and its exact Gram}

At cutoff $N$ retain the \emph{arithmetic}, $S$-monic orthogonal
polynomials $p_n^{\rm ar}(S)$ and their full norms $\omega_n>0$.
As in FW1--5 and HG6--7 \cite{lgamma:FW,lgamma:HG}, use the original classes
\begin{equation}\label{lgamma:eq:nativetriple}
 a_1=b_N=[p_N^{\rm ar}],\qquad
 a_2=b_{N+1}=[p_{N+1}^{\rm ar}],\qquad
 v=v_\lambda=\left[\frac{\chi_k(S)}{(S-\lambda)\chi_k'(\lambda)}\right],
 \quad\lambda=k\rho.
\end{equation}
Brackets denote the original class, with the physical presentation
\eqref{lgamma:eq:unitframe} whenever that is the retained frame.  In particular
the complete physical unit multiplies these classes as specified there.
The vector $a_2$ is a class represented by a degree-$N+1$ polynomial,
but its minimum at cutoff $N$ includes all its degree-at-most-$N$
representatives; they exist because $N\ge q-1$.  It is not replaced
by a Gamma polynomial or by an inadmissible unreduced lift.

Write $G=G_N^{\rm ar}$, $Q=Q_N^{\rm ar}$, and
\begin{equation}\label{lgamma:eq:omegal}
 \Omega=\Lambda^*Q\Lambda,\qquad L=G-\Omega,\qquad
 X=(a_1,a_2,v),\qquad W=\Lambda X.
\end{equation}
The exact class-space minimum section is
$S_B=G^{-1}\Lambda^*Q$.  It satisfies $\Lambda S_B=I_B$.
Thus $P_K=I_E-S_B\Lambda$ maps onto $\ker\Lambda$, is the identity
there, and $I_K^*GS_B=0$ for the actual kernel inclusion
$I_K:\ker\Lambda\hookrightarrow E$.
The resulting orthogonal decomposition proves
\begin{equation}\label{lgamma:eq:split}
 0\preceq\Omega\preceq G,\quad L=GP_K\succeq0,
 \qquad X^*GX=X^*\Omega X+X^*LX.
\end{equation}
This is the exact FW37--38 and HR2 minimum, now proved in the actual
three columns.  Define all its entries without dropping any:
\begin{equation}\label{lgamma:eq:entries}
 \begin{gathered}
 d_i=a_i^*Ga_i,\quad E=v^*Gv,\quad F_i=a_i^*Gv,\\
 D_i=a_i^*\Omega a_i,\quad e=v^*\Omega v,\quad
 B_i=a_i^*\Omega v,\quad z=a_1^*\Omega a_2,\\
 k_i=a_i^*La_i=d_i-D_i,\quad
 K_i=a_i^*Lv=F_i-B_i,\quad e_K=v^*Lv=E-e.
 \end{gathered}
\end{equation}
The native source pairing $a_1^*Ga_2$ is zero, but the observed pairing
$z$ need not be zero.  To verify the former assertion in the original
objects, the exact functional equations $\xi(1-s)=\xi(s)$ and
$h(1-s)=h(s)$ make $w_h$ even, and convolution preserves evenness.
Evenness of $m_k$ makes $p_n^{\rm ar}(k-S)=(-1)^np_n^{\rm ar}(S)$.
Reflection permutes the full roots of $\chi_k$ and takes each entire
affine fibre to the reflected fibre isometrically.  Its induced map
$\Sigma$ therefore satisfies $\Sigma^*G\Sigma=G$, while
$\Sigma b_n=(-1)^nb_n$.  Hence
$a_1^*Ga_2=(-1)^{2N+1}a_1^*Ga_2=0$.  The same assertion in the
physical frame follows by the exact congruence \eqref{lgamma:eq:unitframe}.
Consequently the complete source, observed and kernel triple Grams are
\begin{equation}\label{lgamma:eq:threegrams}
 \begin{gathered}
 \mathcal F=X^*GX=
 \begin{pmatrix}d_1&0&F_1\\0&d_2&F_2\\\bar F_1&\bar F_2&E\end{pmatrix},\\
 \mathcal C=W^*QW=
 \begin{pmatrix}D_1&z&B_1\\\bar z&D_2&B_2\\\bar B_1&\bar B_2&e\end{pmatrix},\qquad
 \mathcal K=X^*LX=
 \begin{pmatrix}k_1&-z&K_1\\-\bar z&k_2&K_2\\\bar K_1&\bar K_2&e_K\end{pmatrix},\\
 \mathcal F=\mathcal C+\mathcal K,\qquad \mathcal C,\mathcal K\succeq0.
 \end{gathered}
\end{equation}
In particular the two complementary off-diagonal boundary terms are
$z$ and $-z$, not two separately vanishing terms.

\begin{lgammatheorem}[The joint native observed constraint]\label{lgamma:thm:joint}
Put
\begin{equation}\label{lgamma:eq:baseline}
 \mathfrak b_{k,N}=\ell_k\beta_N,\qquad
 \Psi_N=W^*\mathcal B_N^{-1}W,\qquad
 \mathcal C^{(0)}=\mathfrak b_{k,N}\Psi_N.
\end{equation}
Then the actual matrices in \eqref{lgamma:eq:threegrams} satisfy
\begin{equation}\label{lgamma:eq:joint}
 \boxed{\mathcal C\succeq\mathcal C^{(0)},\qquad
 \mathcal R:=\mathcal C-\mathcal C^{(0)}\succeq0,\qquad
 \mathcal F-\mathcal C^{(0)}=\mathcal K+\mathcal R,\quad
 \mathcal K\succeq0.}
\end{equation}
All entries of $\Psi_N$ use the same original $\Lambda a_1$,
$\Lambda a_2$, $\Lambda v$ and the complete original map in
\eqref{lgamma:eq:nativelower}.
Moreover $\ker\mathcal C=\ker\Psi_N=\ker W$; no independence of
the three observed columns is assumed.
\end{lgammatheorem}
\begin{proof}
Congruence of \eqref{lgamma:eq:nativelower} by the three-column matrix $W$
gives the first inequality and therefore the positive remainder.
The remaining equality follows from \eqref{lgamma:eq:threegrams}.  It
preserves the whole native kernel remainder rather than replacing it
with a trace or a rank.
Finally, for $t\in\lgammaC^3$, $t^*\mathcal Ct=(Wt)^*Q(Wt)$ vanishes
exactly when $Wt=0$, since $Q>0$.  The same argument with
$\mathcal B_N^{-1}>0$ proves the stated kernel identities.
\end{proof}

Here are all scalar positive-semidefinite constraints in explicit form.
Write $\psi_{ij}=(\Psi_N)_{ij}$ and set
\begin{equation}\label{lgamma:eq:resentries}
 \begin{gathered}
 x=D_1-\mathfrak b_{k,N}\psi_{11},\quad
 y=D_2-\mathfrak b_{k,N}\psi_{22},\quad
 h_B=e-\mathfrak b_{k,N}\psi_{33},\\
 \zeta_B=z-\mathfrak b_{k,N}\psi_{12},\qquad
 r_1=B_1-\mathfrak b_{k,N}\psi_{13},\quad
 r_2=B_2-\mathfrak b_{k,N}\psi_{23}.
 \end{gathered}
\end{equation}
Then $\mathcal R=\left(\begin{smallmatrix}
x&\zeta_B&r_1\\\bar\zeta_B&y&r_2\\\bar r_1&\bar r_2&h_B
\end{smallmatrix}\right)$, and \eqref{lgamma:eq:joint} is equivalent, as far
as its observed lower-form inequality is concerned, to the seven real
conditions
\begin{equation}\label{lgamma:eq:allminors}
 \boxed{\begin{gathered}
 x\ge0,\quad y\ge0,\quad h_B\ge0,\\
 xy-|\zeta_B|^2\ge0,\quad xh_B-|r_1|^2\ge0,
 \quad yh_B-|r_2|^2\ge0,\\
 xyh_B+2\Re(\zeta_Br_2\bar r_1)
 -x|r_2|^2-y|r_1|^2-h_B|\zeta_B|^2\ge0.
 \end{gathered}}
\end{equation}
The last line is the complete determinant, including its triple product.
Necessity follows from principal minors of a positive semidefinite
matrix.  For sufficiency, add $\epsilon I$ with $\epsilon>0$.
Its first and second leading principal minors are positive by the
first two lines.  Its determinant is
\[
 \det(\mathcal R+\epsilon I)
 =\det\mathcal R+
 \epsilon\sum_{|I|=2}\det\mathcal R_{I,I}
 +\epsilon^2\operatorname{tr}\mathcal R+\epsilon^3>0.
\]
Sylvester's criterion gives
$\mathcal R+\epsilon I>0$; let $\epsilon\downarrow0$.
Thus no positivity condition is missing at singular boundary cases.
These are exact constraints on the native matrix; they do not assert
that every numerical matrix satisfying them is realized by the native
arithmetic source.

For the un-subtracted observed Gram the corresponding explicit conditions
are
\begin{equation}\label{lgamma:eq:observedminors}
 \begin{gathered}
 D_1,D_2,e\ge0,\quad D_1D_2\ge|z|^2,\quad
 D_1e\ge|B_1|^2,\quad D_2e\ge|B_2|^2,\\
 D_1D_2e+2\Re(zB_2\bar B_1)
 -D_1|B_2|^2-D_2|B_1|^2-e|z|^2\ge0.
 \end{gathered}
\end{equation}
The kernel remainder is independently retained through
\begin{equation}\label{lgamma:eq:kernelminors}
 \begin{gathered}
 k_1,k_2,e_K\ge0,\quad k_1k_2\ge|z|^2,\quad
 k_1e_K\ge|K_1|^2,\quad k_2e_K\ge|K_2|^2,\\
 k_1k_2e_K-2\Re(zK_2\bar K_1)
 -k_1|K_2|^2-k_2|K_1|^2-e_K|z|^2\ge0.
 \end{gathered}
\end{equation}
The sign change in the last triple product is forced by the exact
off-diagonal $-z$ in \eqref{lgamma:eq:threegrams}.

For $h_B>0$, completing the square in the last coordinate of
$\mathcal R$ gives the equivalent two-dimensional Schur conditions
\begin{equation}\label{lgamma:eq:resdisk}
 x-\frac{|r_1|^2}{h_B}\ge0,\quad
 y-\frac{|r_2|^2}{h_B}\ge0,\qquad
 \left|\bar r_1r_2-h_B\bar\zeta_B\right|^2
 \le(h_Bx-|r_1|^2)(h_By-|r_2|^2).
\end{equation}
No phase was chosen: the disk has the actual complex centre
$h_B\bar\zeta_B$.  If $h_B=0$, \eqref{lgamma:eq:allminors} instead forces
$r_1=r_2=0$ and leaves the exact conditions $x,y\ge0$ and
$xy\ge|\zeta_B|^2$.  This covers the case without a division.
Similarly, when $e>0$ the un-subtracted constraint implies
\begin{equation}\label{lgamma:eq:Bdisk}
 |\bar B_1B_2-e\bar z|^2
 \le(eD_1-|B_1|^2)(eD_2-|B_2|^2).
\end{equation}
Here $e>0$ is equivalent to $\Lambda v\ne0$; when $e=0$, positivity
forces $B_1=B_2=0$.  No additional period locus is silently imposed.

\section{Both complementary cross products and every remainder term}

Retain the native scale $\omega_N>0$ from \eqref{lgamma:eq:nativetriple}.
The exact current formulas in HG7 \cite{lgamma:HG}, also obtained by
substituting HR5 into the original response products \cite{lgamma:HR}, are
\begin{equation}\label{lgamma:eq:currents}
 \begin{aligned}
 \Phi_{K,N}&=\frac2{\omega_N}\Re(\bar K_1K_2),\\
 \Phi_{B,N}&=\frac2{\omega_N}\Re(\bar B_1B_2),\\
 \Phi_{\times,N}&=\frac2{\omega_N}
        \Re(\bar K_1B_2+\bar B_1K_2).
 \end{aligned}
\end{equation}
For verification on the native nonzero-denominator domain, put
$\mathfrak c_N=\bar F_1F_2/(\omega_NE)$, $U=K_1/F_1$ and
$V=K_2/F_2$.  Multiplying $2E\mathfrak c_N$ respectively by
$\bar UV$, $(1-\bar U)(1-V)$, and
$\bar U(1-V)+(1-\bar U)V$ gives exactly
\eqref{lgamma:eq:currents}.  The undivided right sides give their exact
polynomial continuation across vanishing response denominators.

Put, without discarding the positive remainder in \eqref{lgamma:eq:joint},
\[
 B_i^0=\mathfrak b_{k,N}\psi_{i3},\qquad
 K_i^0=F_i-B_i^0,\qquad B_i=B_i^0+r_i,\quad K_i=K_i^0-r_i.
\]
Expansion of every product yields
\begin{equation}\label{lgamma:eq:currentremainders}
 \begin{aligned}
 \frac{\omega_N}2\Phi_{B,N}
 &=\Re\bigl(\bar B_1^0B_2^0+\bar B_1^0r_2
              +\bar r_1B_2^0+\bar r_1r_2\bigr),\\
 \frac{\omega_N}2\Phi_{K,N}
 &=\Re\bigl(\bar K_1^0K_2^0-\bar K_1^0r_2
              -\bar r_1K_2^0+\bar r_1r_2\bigr),\\
 \frac{\omega_N}2\Phi_{\times,N}
 &=\Re\bigl(\bar K_1^0B_2^0+\bar B_1^0K_2^0
       +(\bar K_1^0-\bar B_1^0)r_2\\
 &\hspace{30mm}+\bar r_1(K_2^0-B_2^0)-2\bar r_1r_2\bigr).
 \end{aligned}
\end{equation}
Every linear cross term and every quadratic remainder is displayed.
In their sum these remainder terms cancel algebraically, not by an
estimate, giving the unchanged total
\begin{equation}\label{lgamma:eq:currenttotal}
 \Phi_{K,N}+\Phi_{B,N}+\Phi_{\times,N}
 =\frac2{\omega_N}\Re(\bar F_1F_2)
 =2E\Re\mathfrak c_N.
\end{equation}
The source coefficient is retained in its exact full-source definition,
not replaced by a leading expression omitting its signed phase remainder.
The baseline $\mathcal C^{(0)}$ is a proved lower form, not a declaration
that $\mathcal R$ is small or zero.

For example, \eqref{lgamma:eq:resdisk} supplies the exact finite enclosure
\begin{equation}\label{lgamma:eq:currentdisk}
 \begin{split}
 \Big|\frac{\omega_N}2\Phi_{B,N}
 -\Re(\bar B_1^0B_2^0+\bar B_1^0r_2+\bar r_1B_2^0
                   +h_B\bar\zeta_B)\Big|
 \le\sqrt{(h_Bx-|r_1|^2)(h_By-|r_2|^2)}.
 \end{split}
\end{equation}
For $h_B=0$ the same formula holds because $r_1=r_2=0$.
This is a correlated constraint retaining the actual remainder entries,
not an assignment of unknown phases.  All identities and inequalities
hold at every cutoff in \eqref{lgamma:eq:window}; insertion into the original
current sum keeps its weights $1,2,\ldots,2,1$, and the original
four-endpoint return keeps $q-1,q,2q-1,2q$.

\section{The separate actual Gamma-convolution reference}

For completeness let $b\ge1$ be any integer.  The same recurrence
argument applies to the actual finite-measure convolution $\sigma^{*b}$,
but this is not the native weight in \eqref{lgamma:eq:nativeweight}.  Its exact
mass is $M_{\Gamma,b}=M_\Gamma^b$.  Fubini and \eqref{lgamma:eq:laplace} give
\[
 \int e^{sy}\,d\sigma^{*b}(y)=M_{\Gamma,b}(\cos s)^{-b/2}.
\]
For
\[
 \sum_{n\ge0}p_n^{[b]}(y)\frac{t^n}{n!}
 =(1+t^2)^{-b/4}e^{y\arctan t},
\]
the same two-generating-function calculation proves
$\gamma_n^{[b]}=M_{\Gamma,b}n!(b/2)_n$ and
\[
 p_{n+1}^{[b]}=yp_n^{[b]}-n(n-1+b/2)p_{n-1}^{[b]}.
\]
Write $\alpha_b=b/2$ and $b_n=\sqrt{(n+1)(n+\alpha_b)}$.
For $n\ge1$,
\[
 b_n^2-(1+b_{n-1})^2
 =2n+\alpha_b-1-2\sqrt{n(n+\alpha_b-1)}\ge0,
\]
because the first compared term is positive and its square minus the
second square is $(\alpha_b-1)^2$.  The positive imaginary-point
recurrence and the initial squares $M_{\Gamma,b}^{-1}$ and
$(\alpha_bM_{\Gamma,b})^{-1}$ therefore prove
\begin{equation}\label{lgamma:eq:conv}
 K_N^{\sigma^{*b}}(i,i)
 \le\frac{1+N\max(1,2/b)}{(\sqrt{2\pi})^b}.
\end{equation}
At $b=2$ the recurrence gives equal imaginary-point squares at every
index, and \eqref{lgamma:eq:conv} is an equality.  No assertion in the native
transfer used $m_k=\sigma^{*k}$.

\section{Exact finite verification and provenance}

The companion \texttt{verify\_exact.py} proves by exact rational
calculation the alternating ratios, closed kernel formula and linear
bounds through $N=512$; it checks original Gram norms, inverse traces
and Loewner differences through $N=10$.  It checks symbolic-$c$
coefficient inverses and determinant phases through $N=8$, a complete
complex original-$S$ jet quotient at $N=4$, and the separate
Gamma-convolution recurrence bounds for $1\le k\le12$, $N\le128$.
These finite checks supplement the all-degree proofs, not replace them.
All printed Gram factors retain their exact common scalar
$\sqrt{2\pi}$ symbolically.  The first original kernel values multiplied
by $\sqrt{2\pi}$ are $1,3,9/2,63/10,63/8,77/8,539/48$.
The separate \texttt{verify\_native\_gram\_identities.py} checks the full
Hermitian determinant, complex Schur-disk identity, singular-case
principal-minor criterion, native source--observed--kernel decomposition,
and all three current expansions as exact symbolic polynomial identities.
It assigns no numerical value to a native arithmetic coefficient or
period phase.

The source-use coverage for this derivation is precise: Berg--Szwarc's
original version-one \LaTeX, Section 3, the complete disk-domination
and kernel-eigenvalue lemma proofs and intervening proposition; CAI's
domain, original unit and observation definitions and CAI16--19;
OB's full supplied text, especially OB13 and OB18; HG and HR's full
supplied texts; FW1--7 and FW37--40; and NC's original unit/frame
definitions and NC24--27 distinguishing its separate proper-source
minimum.  No separate full-paper reading of Berg--Szwarc by this
independent derivation is claimed; the parent task read that source in
full.  No Gamma-polynomial theorem from CCM has been imported.

\begin{thebibliography}{9}
\bibitem{lgamma:BS}
Christian Berg and Ryszard Szwarc,
\emph{The smallest eigenvalue of Hankel matrices},
arXiv:0906.4506v1 (2009), Section 3, original-source labels
\texttt{thm:pni} and \texttt{thm:lambda}, both lemmas;
\url{https://arxiv.org/abs/0906.4506v1}.
The original author \LaTeX\ source is retained locally as
\path{sources/0906.4506v1/original_source/main.tex}.

\bibitem{lgamma:CAI}
Split-Zero research programme,
\emph{The complete original action relative to its arithmetic total},
complete native CAI provider reconstructed from the delivered continuation
of 16 September 2026, CAI1--3, CAI6--9 and CAI16--19.
Original provider: \path{providers/CAI_COMPLETE_PREREQUISITE.tex}.
The retained analytic constants are the source's TW7 and BT12 constants,
as recorded in its AMT4/CAI16--17.

\bibitem{lgamma:OB}
Split-Zero research programme,
\emph{The observed terminal class on the complete original window},
native provider \path{OBSERVED_CLASS_RETURN.tex}, OB1--2,
the degree-$2q+2$ comparison preceding OB13, and OB18--19.

\bibitem{lgamma:FW}
Split-Zero research programme,
\emph{Full-window source products and the retained class:
finite quotient comparisons in the original arithmetic metric},
19 September 2026, FW1--7 and FW37--40.
Original provider: \path{providers/FULL_WINDOW_SOURCE_PRODUCTS.tex}.

\bibitem{lgamma:HG}
Split-Zero research programme,
\emph{Subexponential heat cost and the original arithmetic currents},
19 September 2026, HG1 and HG6--11.
Original provider: \path{providers/HEAT_GROWTH_AND_ORIGINAL_CURRENTS.tex}.

\bibitem{lgamma:HR}
Split-Zero research programme,
\emph{Heat approximation errors for the original complex observation responses},
19 September 2026, HR2--5 and HR8--11.
Original provider: \path{providers/HEAT_OBSERVATION_RESPONSE_ERRORS.tex}.

\bibitem{lgamma:NC}
Split-Zero research programme,
\emph{Vanishing of the original proper-source kernel through conductor coprimality},
19 September 2026, NC1--3 and NC24--27.
Original provider: \path{providers/ORIGINAL_CONDUCTOR_COPRIMALITY.tex}.
Used here only to retain the actual unit/frame conventions and to avoid
substituting its separate proper-source minimum for the present fibre.
\end{thebibliography}

\endgroup

\clearpage\begingroup
\part*{Closed range of the original full theta source}
\newtheorem{lthetatheorem}{Theorem}
\newtheorem{lthetalemma}[lthetatheorem]{Lemma}
\newtheorem{lthetaproposition}[lthetatheorem]{Proposition}
\newtheorem{lthetacorollary}[lthetatheorem]{Corollary}
\newcommand{\lthetaB}{\mathscr B}
\newcommand{\lthetaSch}{\mathcal S}
\newcommand{\lthetaHH}{\mathcal H}
\newcommand{\lthetaZ}{\mathbb Z}
\newcommand{\lthetaR}{\mathbb R}
\newcommand{\lthetaC}{\mathbb C}
\newcommand{\lthetaN}{\mathbb N}
\newcommand{\lthetaid}{\operatorname{id}}
\newcommand{\lthetacl}{\operatorname{cl}}
\newcommand{\lthetasupp}{\operatorname{supp}}
\newcommand{\lthetaF}{\mathcal F}



\section{The fixed objects and the source theorem}

Every space and map in this note is over $\lthetaC$. The original source is
\[
 V=\left\{\phi\in\lthetaSch(\lthetaR):\phi(-x)=\phi(x),\quad
       \phi(0)=0,\quad \int_\lthetaR\phi(x)\,dx=0\right\},
\]
with the subspace topology from the ordinary Schwartz space. Put
\[
 D=-x\frac{d}{dx},\qquad
 p_{b,j}(F)=\sup_{x>0}x^b|D^jF(x)|\quad(b\in\lthetaZ,\ j\in\lthetaN_0),
\]
\[
 \lthetaB=\{F\in C^\infty(\lthetaR_{>0}):p_{b,j}(F)<\infty
                     \text{ for every }b,j\}.
\]
The topology on $\lthetaB$ is exactly that generated by all $p_{b,j}$. The
operators and Fourier convention are
\begin{equation}\label{ltheta:eq:original}
 \Theta\phi(x)=2\sum_{n=1}^{\infty}\phi(nx),\qquad
 JF(x)=x^{-1}F(x^{-1}),\qquad
 \widehat\phi(\xi)=\int_\lthetaR\phi(x)e^{-2\pi i x\xi}\,dx.
\end{equation}
No conjugation by $x^{1/2}$, renormalization of $\Theta$, change of the
coordinate $x$, or substitution of a Hilbert norm is made here.

The literature source is Ralf Meyer, \emph{A spectral interpretation for
the zeros of the Riemann zeta function}, arXiv:math/0412277v3,
\href{https://arxiv.org/abs/math/0412277v3}{original versioned source}.
Its Section 2 defines
\[
 \lthetaHH_- =\bigcap_{s\in\lthetaR}\lthetaSch(\lthetaR_+^*)_s,\qquad
 \lthetaSch(\lthetaR_+^*)_s=
 \{F:x\mapsto x^sF(x)\text{ is Schwartz in }\log x\}.
\]
Its Section 3 defines $Z\phi(x)=\sum_{n\geq1}\phi(nx)$ and
$\lthetaHH_\cap=\{\phi\in\lthetaSch(\lthetaR):\phi\text{ even},\,
\phi(0)=\lthetaF\phi(0)=0\}$. Theorem 3.3, with source label
\texttt{the:Zeta\_estimate}, proves closed range and the
topological-isomorphism-onto-range property, including the restriction
$Z:\lthetaHH_\cap\to\lthetaHH_-$. Its proof uses Poisson summation and
Proposition 3.2 (label \texttt{pro:Z\_invertible}), together with
Proposition 2.1 (label \texttt{pro:Sch\_estimates}).

Meyer's Fourier kernel has the opposite sign. For an even $\phi$ the
change of variable $x\mapsto-x$ in the defining integral proves that
the two transforms are equal, as functions, with no scalar factor.
Consequently $V=\lthetaHH_\cap$ with identical source topology. His $Z$ and
the original map satisfy $\Theta=2Z$. The next sections prove the target
topology identification and independently reconstruct the closed-range
conclusion with explicit inverse estimates.

\section{Identity equivalence of the complete seminorm systems}

For $s\in\lthetaR$ and $k,j\in\lthetaN_0$ define
\begin{equation}\label{ltheta:eq:q}
 q_{s,k,j}(F)=\sup_{x>0}(1+|\log x|)^k
 \left|\left(x\frac d{dx}\right)^j(x^sF(x))\right|.
\end{equation}
These are the ordinary Schwartz seminorms in the variable $\log x$
after multiplication by $x^s$, expressed on the original $x$-axis. Their
projective family generates precisely Meyer's topology on $\lthetaHH_-$.

\begin{lthetatheorem}\label{ltheta:thm:topology}
The identity on positive-axis functions is a topological linear
isomorphism $\lthetaB=\lthetaHH_-$. More explicitly, set
\[
 n_-(s)=\lfloor s\rfloor-1,\qquad n_+(s)=\lceil s\rceil+1,
 \qquad C_k=\sup_{t\geq0}(1+t)^ke^{-t}.
\]
Here $C_0=1$ and $C_k=k^ke^{1-k}$ for integer $k\geq1$.
For every real $s$ and every $k,j\geq0$,
\begin{equation}\label{ltheta:eq:q-from-p}
 q_{s,k,j}(F)\leq C_k\sum_{\ell=0}^j\binom j\ell
 |s|^{j-\ell}\bigl(p_{n_-(s),\ell}(F)+p_{n_+(s),\ell}(F)\bigr).
\end{equation}
Conversely, for every integer $b$ and $j\geq0$,
\begin{equation}\label{ltheta:eq:p-from-q}
 p_{b,j}(F)\leq\sum_{\ell=0}^j\binom j\ell
 |b|^{j-\ell}q_{b,0,\ell}(F).
\end{equation}
In these formulas a zeroth power, including $0^0$, means $1$.
\end{lthetatheorem}

\begin{proof}
Write $L=x\,d/dx=-D$. The Leibniz rule, iterated $j$ times, gives the
exact operator identity
\[
 L^j(x^sF)=x^s(s-D)^jF
   =x^s\sum_{\ell=0}^j\binom j\ell s^{j-\ell}(-1)^\ell D^\ell F.
\]
For $x\geq1$ put $t=\log x$. Since $n_+(s)-s\geq1$,
\[
 (1+|\log x|)^kx^s\leq C_kx^{n_+(s)}.
\]
For $0<x\leq1$ put $t=-\log x$. Since $s-n_-(s)\geq1$,
\[
 (1+|\log x|)^kx^s\leq C_kx^{n_-(s)}.
\]
Substitution in the displayed binomial expansion proves
\eqref{ltheta:eq:q-from-p}. The constant follows by maximizing
$k\log(1+t)-t$ on $t\geq0$.

The conjugate identity $(b-L)(x^bF)=x^bDF$ yields
\[
 x^bD^jF=(b-L)^j(x^bF).
\]
Its binomial expansion gives \eqref{ltheta:eq:p-from-q}. Hence finiteness of
the full $p$-family and of the full $q$-family is equivalent, and each
seminorm of either family is bounded by a finite linear combination
from the other. This proves equality of the sets and of their locally
convex topologies, through the identity map.
\end{proof}

For completeness, $\lthetaB$ is a Fr\'echet space directly in the $p$-topology.
The seminorm family is countable and separates points. If $(F_n)$ is
Cauchy in every $p_{b,j}$, it is Cauchy in $C^\infty(K)$ for each compact
$K\subset(0,\infty)$, because
\begin{equation}\label{ltheta:eq:R}
 x^m\partial_x^m=R_m(D),\qquad
 R_m(t)=(-1)^m\prod_{r=0}^{m-1}(t+r),\qquad R_0(t)=1.
\end{equation}
The identity follows by induction, using
$x^{m+1}\partial^{m+1}=(x\partial-m)x^m\partial^m$.
Local limits therefore give a single smooth $F$. For a fixed $b,j$,
the $p_{b,j}$-Cauchy bound on $F_n-F_m$ passes pointwise to $m\to\infty$
and then to the supremum; thus $p_{b,j}(F_n-F)\to0$ and $F\in\lthetaB$.
The same local-limit argument proves completeness of $\lthetaSch(\lthetaR)$.
Parity, evaluation at zero, and integration are continuous Schwartz
functionals/conditions (for integration use the seminorm with weight
$(1+|x|)^2$). Therefore $V$ is a closed, complete Fr\'echet subspace.

\section{Continuity, Poisson symmetry, and exact M\"obius inversion}

Write
\[
 T_{N,j}(\psi)=\sup_{t\geq1}t^N|D^j\psi(t)|
 \qquad(\psi\in\lthetaSch(\lthetaR)).
\]
Each $T_{N,j}$ is a continuous Schwartz seminorm: expanding
$D^j=(-x\partial_x)^j$ gives a finite sum of polynomial coefficients
times ordinary derivatives. For $N>1$, termwise differentiation of
the theta series on every compact subset of $(0,\infty)$ is justified
by the convergent majorant $C n^{-N}$. On $x\geq1$ it gives
\begin{equation}\label{ltheta:eq:theta-tail}
 |D^j\Theta\psi(x)|\leq2\zeta(N)x^{-N}T_{N,j}(\psi).
\end{equation}

The ordinary Fourier transform maps $\lthetaSch(\lthetaR)$ continuously to itself.
Indeed, its differentiated, polynomially weighted expressions are
Fourier transforms of finite sums of derivatives of polynomially
weighted $\psi$, multiplied by the explicit powers of $2\pi i$ from
the kernel in \eqref{ltheta:eq:original}; their suprema are bounded by the
corresponding $L^1$ norms. Those norms are bounded by finitely many
Schwartz seminorms, using the integrable weight $(1+|x|)^{-2}$.
Fourier inversion gives $\widehat{\widehat\psi}(x)=\psi(-x)$.
It follows that the Fourier transform is an involution of $V$:
evenness is preserved, $\widehat\phi(0)=\int\phi=0$, and
$\int\widehat\phi=\phi(0)=0$.

Poisson summation with the exact kernel \eqref{ltheta:eq:original} says
\[
 \sum_{n\in\lthetaZ}\psi(nx)
    =x^{-1}\sum_{n\in\lthetaZ}\widehat\psi(n/x).
\]
It follows, for example, by periodizing $t\mapsto\psi(xt)$:
its $k$th Fourier-series coefficient is $x^{-1}\widehat\psi(k/x)$,
and the series and all needed derivatives converge absolutely by
Schwartz decay; evaluation at $t=0$ gives the formula.
For even $\psi$ this is exactly
\begin{equation}\label{ltheta:eq:poisson}
 \psi(0)+\Theta\psi(x)
       =x^{-1}\widehat\psi(0)+J\Theta\widehat\psi(x).
\end{equation}
Both endpoint terms vanish for $\phi\in V$, so
\begin{equation}\label{ltheta:eq:symmetry}
 \Theta\phi=J\Theta\widehat\phi,
 \qquad J\Theta\phi=\Theta\widehat\phi.
\end{equation}
The second identity uses $J^2=\lthetaid$.

Direct differentiation on the original axis gives
\begin{equation}\label{ltheta:eq:J}
 D^jJ=J(1-D)^j,\qquad
 p_{b,j}(JF)\leq\sum_{r=0}^j\binom jr p_{1-b,r}(F).
\end{equation}
For $j=1$ the first identity follows from
$-x\partial_x(x^{-1}F(x^{-1}))
=x^{-1}F(x^{-1})+x^{-2}F'(x^{-1})$;
induction proves the general case. In the seminorm bound substitute
$y=x^{-1}$, so $x^b x^{-1}=y^{1-b}$, and expand $(1-D)^j$.
In particular $J$ is a continuous involution of $\lthetaB$.

Choose the integers $N=\max(2,b)$ and $N'=\max(2,1-b)$.
Combining \eqref{ltheta:eq:theta-tail} on $x\geq1$ with
\eqref{ltheta:eq:symmetry} and \eqref{ltheta:eq:J} on $0<x\leq1$ yields
\begin{equation}\label{ltheta:eq:theta-continuity}
 p_{b,j}(\Theta\phi)
 \leq2\zeta(N)T_{N,j}(\phi)
  +2\zeta(N')\sum_{r=0}^j\binom jrT_{N',r}(\widehat\phi).
\end{equation}
Thus $\Theta:V\to\lthetaB$ is well defined and continuous in the precise
original topologies.

For a positive-axis smooth function define
$p^+_{u,j}(\psi)=\sup_{x>0}x^u|D^j\psi(x)|$ whenever finite;
the superscript distinguishes a source estimate from the target
seminorm notation. Let $\mu(n)$ be the M\"obius function and put
\[
 (MG)(x)=\sum_{n=1}^{\infty}\mu(n)G(nx).
\]
Whenever $p^+_{u,r}(G)<\infty$ for $u>1$ and $0\leq r\leq j$, its
series may be differentiated $j$ times locally uniformly and
\begin{equation}\label{ltheta:eq:mobius-bound}
 p^+_{u,j}(MG)\leq\zeta(u)p^+_{u,j}(G).
\end{equation}
Here $|\mu(n)|\leq1$ and $D^j[G(nx)]=(D^jG)(nx)$ give the estimate
term by term. In particular this applies to every $G\in\lthetaB$ and
integer $u\geq2$.

For $\phi\in V$ and fixed $x>0$, the double series
\[
 \sum_{n,m\geq1}|\mu(n)\phi(mnx)|
 \leq x^{-u}p^+_{u,0}(\phi)\zeta(u)^2
 \qquad(u>1)
\]
is finite. Grouping by $k=mn$ and using
$\sum_{n\mid k}\mu(n)=1$ for $k=1$ and $0$ otherwise gives
\begin{equation}\label{ltheta:eq:mobius-inversion}
 M\Theta\phi=2\phi\quad\text{on }(0,\infty).
\end{equation}
The divisor identity follows by taking the product
$\prod_{p\mid k}(1-1)$ when $k>1$; for $k=1$ the empty product is $1$.
Thus $\Theta$ is injective. Writing $F=\Theta\phi$, the exact inverse
formulas and their estimates are
\begin{align}
 \phi(x)&=\tfrac12\sum_{n\geq1}\mu(n)F(nx),&
 p^+_{u,j}(\phi)&\leq\tfrac12\zeta(u)p_{u,j}(F),
       \label{ltheta:eq:inverse-phi}\\
 \widehat\phi(x)&=\tfrac12\sum_{n\geq1}\mu(n)(JF)(nx),&
 p^+_{u,j}(\widehat\phi)&\leq\tfrac12\zeta(u)
                    \sum_{\ell=0}^j\binom j\ell p_{1-u,\ell}(F),
       \label{ltheta:eq:inverse-hat}
\end{align}
for $x>0$, integers $u\geq2$, and $j\geq0$. The second line uses
\eqref{ltheta:eq:symmetry}. These are inverse formulas on the image; they
do not assert that $MG$ extends to an element of $V$ for an arbitrary
$G\in\lthetaB$.

\section{Explicit recovery of every source Schwartz seminorm}

The following estimate supplies the part of the inverse-continuity
argument where both the function and its Fourier transform are needed.
All $L^2$ norms in this section use ordinary Lebesgue measure $dx$ or
$d\xi$ on $\lthetaR$, and the Fourier transform is exactly
\eqref{ltheta:eq:original}. We use its Plancherel norm equality.

\begin{lthetalemma}\label{ltheta:lem:uncertainty}
For $h\in\lthetaSch(\lthetaR)$ and integers $M,N\geq0$,
\begin{equation}\label{ltheta:eq:uncertainty}
 \|h\|_2\leq2\left(4^M\|x^Mh\|_2
                         +4^N\|\xi^N\widehat h\|_2\right).
\end{equation}
\end{lthetalemma}

\begin{proof}
Let $P$ and $Q$ be multiplication by the indicator of $[-1/4,1/4]$
on the physical and Fourier sides, respectively. The operator
$P\lthetaF^{-1}Q$ has kernel $1_{[-1/4,1/4]}(x)
e^{2\pi i x\xi}1_{[-1/4,1/4]}(\xi)$; its Hilbert--Schmidt norm is
$\sqrt{(1/2)(1/2)}=1/2$. Therefore, by Plancherel and the triangle
inequality,
\[
 \|h\|_2\leq\|Ph\|_2+\|(1-P)h\|_2
 \leq\tfrac12\|h\|_2+\|(1-Q)\widehat h\|_2+\|(1-P)h\|_2.
\]
On $|x|>1/4$, $1\leq4^M|x|^M$; the analogous statement holds in
frequency. Moving the half-norm to the left proves the assertion.
\end{proof}

Write the exact polynomial from \eqref{ltheta:eq:R} as
$R_m(t)=\sum_{j=0}^m r_{m,j}t^j$. No coefficient is dropped in what
follows. For integers $c\geq3$ and $m\geq0$ put
\begin{align}
 E_{c,m}(F)&=\frac1{\sqrt2}\sum_{j=0}^m|r_{m,j}|
 \left(\frac{\zeta(c-1)}{\sqrt3}p_{c-1,j}(F)
                      +\zeta(c+1)p_{c+1,j}(F)\right),
                     \label{ltheta:eq:E}\\
 E^J_{c,m}(F)&=\frac1{\sqrt2}\sum_{j=0}^m|r_{m,j}|
 \left(\frac{\zeta(c-1)}{\sqrt3}
                     \sum_{\ell=0}^j\binom j\ell p_{2-c,\ell}(F)
       +\zeta(c+1)\sum_{\ell=0}^j\binom j\ell p_{-c,\ell}(F)\right).
                     \label{ltheta:eq:EJ}
\end{align}
For $a,b\in\lthetaN_0$ set, with $(b)_r=b!/(b-r)!$,
\begin{multline}\label{ltheta:eq:A}
 A_{a,b}(F)=2\left[4^{b+3}E_{a+3,b}(F)\right.\\
 \left.{}+4^{a+3}(2\pi)^{b-a}
 \sum_{r=0}^{\min(a,b)}\binom ar(b)_rE^J_{b+3,a-r}(F)\right].
\end{multline}
Every $A_{a,b}$ is explicitly a finite nonnegative linear combination
of original $p_{b,j}$ seminorms.

\begin{lthetaproposition}\label{ltheta:prop:inverse}
For $\phi\in V$, $F=\Theta\phi$, and $a,b\in\lthetaN_0$,
\begin{equation}\label{ltheta:eq:L2-inverse}
 \|x^a\phi^{(b)}\|_2\leq A_{a,b}(F).
\end{equation}
In particular the ordinary Schwartz seminorms satisfy
\begin{equation}\label{ltheta:eq:sup-inverse}
 \sup_{x\in\lthetaR}|x^a\phi^{(b)}(x)|
 \leq A_{a,b}(F)+aA_{a-1,b}(F)+A_{a,b+1}(F),
\end{equation}
where the middle term is omitted when $a=0$.
\end{lthetaproposition}

\begin{proof}
For any even Schwartz $\psi$, splitting the positive-axis integral at
$1$ gives
\[
 \|x^cD^j\psi\|_{L^2(\lthetaR)}
 \leq\sqrt2\left(\frac{p^+_{c-1,j}(\psi)}{\sqrt3}
                               +p^+_{c+1,j}(\psi)\right).
\]
Indeed, on $(0,1)$ the integrand's absolute value is at most
$xp^+_{c-1,j}$, whose squared integral is $(p^+_{c-1,j})^2/3$;
on $(1,\infty)$ it is at most $x^{-1}p^+_{c+1,j}$, whose squared
integral is $(p^+_{c+1,j})^2$. Evenness supplies the factor $\sqrt2$;
the square root of the sum is at most the sum of square roots.
Using the full polynomial $R_m$ and then
\eqref{ltheta:eq:inverse-phi}--\eqref{ltheta:eq:inverse-hat}, with $c-1\geq2$, proves
\begin{equation}\label{ltheta:eq:U-bound}
 \|x^cR_m(D)\phi\|_2\leq E_{c,m}(F),\qquad
 \|x^cR_m(D)\widehat\phi\|_2\leq E^J_{c,m}(F).
\end{equation}

Apply Lemma~\ref{ltheta:lem:uncertainty} to
$h=x^a\phi^{(b)}$, with $M=b+3$ and $N=a+3$. The first weighted
term is, by \eqref{ltheta:eq:R},
\[
 x^Mh=x^{a+b+3}\phi^{(b)}=x^{a+3}R_b(D)\phi.
\]
The exact Fourier identities for the specified negative-sign kernel are
\[
 \widehat{\partial_x^b\phi}(\xi)=(2\pi i\xi)^b\widehat\phi(\xi),
 \qquad
 \widehat{x^a\psi}(\xi)=\left(-\frac1{2\pi i}\right)^a
                                      \partial_\xi^a\widehat\psi(\xi).
\]
Their Leibniz expansion therefore gives
\[
 \xi^{a+3}\widehat h(\xi)
 =\left(-\frac1{2\pi i}\right)^a(2\pi i)^b
 \sum_{r=0}^{\min(a,b)}\binom ar(b)_r
       \xi^{a+3+b-r}\widehat\phi^{(a-r)}(\xi).
\]
For each summand the final factor is exactly
$\xi^{b+3}R_{a-r}(D)\widehat\phi(\xi)$, where
$D=-\xi\partial_\xi$ on this axis. The scalar prefactor has modulus
$(2\pi)^{b-a}$. Substituting \eqref{ltheta:eq:U-bound} into
\eqref{ltheta:eq:uncertainty} proves \eqref{ltheta:eq:L2-inverse} with exactly
\eqref{ltheta:eq:A}.

For a Schwartz function $h$, integration from $-\infty$ to $x$ gives
$|h(x)|^2\leq2\|h\|_2\|h'\|_2$, whence
$\|h\|_\infty\leq\|h\|_2+\|h'\|_2$.
For $h=x^a\phi^{(b)}$ its derivative is
$a x^{a-1}\phi^{(b)}+x^a\phi^{(b+1)}$. The preceding $L^2$ estimate
and the triangle inequality yield \eqref{ltheta:eq:sup-inverse}. The
seminorms on its left generate the full Schwartz topology; weights
$(1+|x|)^a$ are bounded by the finite sum
$\sum_{r=0}^a\binom ar|x|^r$ and conversely each $|x|^a$ is bounded by
$(1+|x|)^a$.
\end{proof}

\section{Closed image, the exact quotient, and the chain isomorphism}

\begin{lthetatheorem}\label{ltheta:thm:closed}
The original map $\Theta:V\to\lthetaB$ is a continuous linear injection,
a topological isomorphism onto its range with the subspace topology,
and its range is closed. In particular
\begin{equation}\label{ltheta:eq:closed}
 \overline{\Theta V}^{\,\lthetaB}=\Theta V.
\end{equation}
\end{lthetatheorem}

\begin{proof}
Continuity is \eqref{ltheta:eq:theta-continuity}, injectivity is
\eqref{ltheta:eq:mobius-inversion}, and inverse continuity is
\eqref{ltheta:eq:sup-inverse}. Suppose $F_n=\Theta\phi_n$ converges in $\lthetaB$
to $F$. Applying \eqref{ltheta:eq:sup-inverse} to
$\phi_n-\phi_m\in V$ shows that $(\phi_n)$ is Cauchy for every
Schwartz seminorm. Completeness and closedness of $V$ give
$\phi_n\to\phi\in V$. Continuity of $\Theta$ implies
$\Theta\phi_n\to\Theta\phi$ in $\lthetaB$, and its Hausdorff property gives
$F=\Theta\phi$. Thus the range is sequentially closed. Since $\lthetaB$ is
metrizable by its countable seminorm family, it is closed.
\end{proof}

The literature application reaches the same conclusion immediately
after Theorem~\ref{ltheta:thm:topology}: Meyer's restricted closed embedding
is precisely $Z=\Theta/2$ with source $V=\lthetaHH_\cap$ and target
$\lthetaHH_-=\lthetaB$. Since $2V=V$, its image equals $\Theta V$ as a subset of
that identical target. The independent estimates above additionally
make the inverse continuity completely explicit in original seminorms.

Let $Q=\lthetaB/\Theta V$ with its quotient topology. It is Hausdorff:
if a coset is nonzero, its representative lies outside the closed
subspace $\Theta V$, and the quotient separates it from the zero
coset. It is also Fr\'echet. One direct completeness argument chooses
an increasing defining seminorm sequence for $\lthetaB$, uses the quotient
seminorms, selects from a quotient-Cauchy sequence a subsequence with
successive quotient distances below $2^{-n}$ in each of the first
$n$ seminorms, and lifts each successive difference to a vector whose
same finitely many seminorms are below $2^{1-n}$. The lifted series is
Cauchy in each seminorm, hence converges in $\lthetaB$ and projects to the
subsequence limit; the original Cauchy sequence has the same limit.
The quotient is metrizable by these countably many quotient seminorms.

The original algebraic exact sequence labelled (E5),
\[
 0\longrightarrow
 \overline{\Theta V}^{\,\lthetaB}/\Theta V\longrightarrow Q
 \longrightarrow\lthetaB/\overline{\Theta V}^{\,\lthetaB}\longrightarrow0,
\]
therefore has its first nontrivial term equal to zero. Its last map
is exactly
\begin{equation}\label{ltheta:eq:E5}
 Q\xrightarrow{\ \cong\ }\lthetaB/\overline{\Theta V}^{\,\lthetaB},
 \qquad [F]_{\Theta V}\longmapsto[F]_{\overline{\Theta V}^{\,\lthetaB}},
\end{equation}
with inverse the same representative map. Both maps are continuous
because the two equivalence relations and quotient topologies are
identical. No new completion or representative change is involved.

To retain the operator factor at the chain level, define two-term
complexes in degrees $0,1$ by
\[
 C_\Theta=[V\xrightarrow{\Theta}\lthetaB],\qquad
 C_Z=[\lthetaHH_\cap\xrightarrow{Z}\lthetaHH_-].
\]
The explicitly typed chain isomorphism is
\begin{equation}\label{ltheta:eq:chain}
 T^0=2\lthetaid:V\longrightarrow\lthetaHH_\cap,\qquad
 T^1=\lthetaid:\lthetaB\longrightarrow\lthetaHH_-.
\end{equation}
Indeed $ZT^0\phi=Z(2\phi)=\Theta\phi=T^1\Theta\phi$.
Its inverse has degree-zero map $\tfrac12\lthetaid$ and degree-one map
$\lthetaid$. Thus the induced degree-one quotient isomorphism is
$[F]\mapsto[F]$, and degree-zero cohomology is zero on both sides
by injectivity. These are the range complexes associated to the
operators; no different complex is attributed to Meyer.

For the original two-leg differential the equally literal diagram is
\[
 [V\oplus V\xrightarrow{\,\Theta-J\Theta\,}\lthetaB]
 \xrightarrow{\ (2\lthetaid\oplus2\lthetaid,\,\lthetaid)\ }
 [\lthetaHH_\cap\oplus\lthetaHH_\cap\xrightarrow{\,Z-JZ\,}\lthetaHH_-].
\]
The equality of differentials follows separately for each signed leg.
Its inverse scales both source coordinates by $1/2$ and fixes the
target. Since $J\Theta=\Theta\lthetaF$ on $V$, the source map
$(\phi_+,\phi_-)\mapsto\phi_+-\widehat\phi_-$ and the target identity
define a surjective chain map from the original two-leg complex to
$C_\Theta$. Its kernel is the degree-zero graph
$\{(\widehat\psi,\psi):\psi\in V\}$ with zero degree-one term:
both inclusions follow directly from
$\phi_+-\widehat\phi_-=0$. In particular the two-leg degree-one image
is also precisely $\Theta V$, not a different closure.

Finally $D(V)\subset V$: it preserves evenness, $D\phi(0)=0$, and
integration by parts gives $\int_\lthetaR D\phi=\int_\lthetaR\phi=0$.
It acts continuously on $\lthetaB$ because $p_{b,j}(DF)=p_{b,j+1}(F)$,
and termwise differentiation gives $D\Theta=\Theta D$. Thus the
one-leg chain isomorphism commutes with $D$ in both degrees.
The signs for the two-leg complex are also explicit. Integration by
parts and differentiation of the Fourier kernel give
\[
 \lthetaF D\phi=(1+\xi\partial_\xi)\widehat\phi
           =(1-D)\lthetaF\phi,
 \qquad \lthetaF(1-D)\phi=D\lthetaF\phi.
\]
Consequently its source generator is $D\oplus(1-D)$ and its target
generator is $D$: by \eqref{ltheta:eq:J},
\[
 D(\Theta\phi_+-J\Theta\phi_-)
    =\Theta D\phi_+-J\Theta(1-D)\phi_-.
\]
The two-leg scaling isomorphism commutes with these generators.
So does the projection to the one-leg complex, because
$\phi_+-\widehat\phi_-$ is sent under the source generator to
$D\phi_+-\lthetaF(1-D)\phi_-=D(\phi_+-\widehat\phi_-)$.
No common diagonal generator on the two source legs and no sign
change to $x\partial_x$ is silently substituted.

\section{The exact scope for subspaces, finite cutoffs, and Hilbert norms}

\begin{lthetaproposition}\label{ltheta:prop:subspace}
For every linear subspace $W\subseteq V$, with closure taken in the
indicated original spaces,
\begin{equation}\label{ltheta:eq:W}
 \overline{\Theta W}^{\,\lthetaB}
       =\Theta\left(\overline W^{\,V}\right).
\end{equation}
The restriction of $\Theta$ induces a topological linear isomorphism
\[
 \overline W^{\,V}/W
       \longrightarrow\overline{\Theta W}^{\,\lthetaB}/\Theta W,
 \qquad[\psi]\longmapsto[\Theta\psi],
\]
with the respective quotient topologies, whether or not the quotients
are Hausdorff. In particular every finite-dimensional $W$ has closed
theta image in $\lthetaB$.
\end{lthetaproposition}

\begin{proof}
The range $R=\Theta V$ is closed in $\lthetaB$ and $\Theta:V\to R$ is a
homeomorphism. Therefore the closure of $\Theta W$ in $\lthetaB$ lies in
$R$ and equals its closure in $R$. A homeomorphism carries closures
to closures, proving \eqref{ltheta:eq:W}. Its restrictions to the closed
subspaces in that equality are inverse homeomorphisms. They carry
the equivalence classes modulo $W$ and $\Theta W$ to each other,
so the quotient universal property proves both induced maps
continuous and inverse.

For the last assertion, finite-dimensional subspaces of the present
Hausdorff locally convex space are closed. Here is a direct proof.
For a basis $w_1,\ldots,w_d$ of $W$, the map
$T:\lthetaC^d\to V$, $c\mapsto\sum c_iw_i$, is continuous and injective.
For each vector on the Euclidean unit sphere choose a continuous
seminorm positive on its image. Compactness gives finitely many
such seminorms whose sum $p$ is positive on every image on the
sphere; continuity and compactness give $p(Tc)\geq\delta|c|$ for
some $\delta>0$. Thus a convergent sequence $Tc_n$ in $V$ has
Cauchy coordinates, with limit $c$, and its limit is $Tc\in W$.
Metrizability gives closedness of $W$. Now apply \eqref{ltheta:eq:W}.
\end{proof}

Consequently for any programme subspace denoted $W_h$, formula
\eqref{ltheta:eq:W}, not an unproved replacement of $W_h$ by $V$, is the
exact available assertion. It computes its closure defect as the
source defect $\overline{W_h}^{\,V}/W_h$. For each fixed finite
polynomial source span $W_L$, its image is closed in $\lthetaB$ by the
finite-dimensional assertion. For a union of these spans the precise
formula remains
\[
 \overline{\bigcup_L\Theta W_L}^{\,\lthetaB}
    =\Theta\left(\overline{\bigcup_L W_L}^{\,V}\right)
\]
when the spans are increasing, so their union is a linear subspace.
No assertion that this source closure is all of $V$ is made.

The inclusion into the original Hilbert space
$\lthetaHH=L^2(\lthetaR_{>0},dx)$ is a different, explicitly continuous map:
\begin{equation}\label{ltheta:eq:Hilbert}
 \|F\|_{\lthetaHH}^2
 \leq p_{0,0}(F)^2\int_0^1dx
          +p_{1,0}(F)^2\int_1^\infty x^{-2}\,dx
 =p_{0,0}(F)^2+p_{1,0}(F)^2.
\end{equation}
It is injective since a continuous function zero almost everywhere
is zero everywhere. Its inverse on its image is not continuous.
Indeed fix $0\ne\eta\in C_c^\infty((1,2))$ and put
$F_n(x)=n^{-1}\eta(x-n)$. Each $F_n$ belongs to $\lthetaB$;
$\|F_n\|_{\lthetaHH}=n^{-1}\|\eta\|_2\to0$, whereas for a fixed
$t_0\in(1,2)$ with $\eta(t_0)\ne0$,
\[
 p_{2,0}(F_n)\geq n^{-1}(n+t_0)^2|\eta(t_0)|\longrightarrow\infty.
\]
Thus Hilbert convergence does not imply convergence in the original
target topology. In particular the closed-image theorem does not
assert Hilbert closedness, alter any finite Hilbert Gram matrix, or
replace a calculation of Hilbert closure. It removes exactly the
closure defect in (E5) for the original full source $V$ and the
original target seminorm topology.


\endgroup

\clearpage\begingroup
\part*{The actual arithmetic Jacobi commutant and its finite boundary}
\newtheorem{ljacobitheorem}{Theorem}
\newtheorem{ljacobiproposition}[ljacobitheorem]{Proposition}
\newcommand{\ljacobiR}{\mathbb R}
\newcommand{\ljacobiC}{\mathbb C}
\newcommand{\ljacobiN}{\mathbb N}
\newcommand{\ljacobiJ}{\mathcal J}
\newcommand{\ljacobiCC}{\mathcal C}
\newcommand{\ljacobiDom}{\operatorname{Dom}}



\section{Unchanged data and the proved envelope}

Retain the actual original packet divisor
\[
 \rho=\tfrac12+\delta+i\gamma,\quad 0<\delta<\tfrac12,\quad\gamma>2,
 \quad h(s)=\prod_{z\in\{\rho,\bar\rho,1-\rho,1-\bar\rho\}}(s-z)^m,
\]
with its full multiplicity $m\geq1$, $g=2\xi$, and
\[
 w_h(y)=\frac{|(g/h)(1/2+iy)|^2}{2\pi},\qquad
 m_k=w_h^{*k},\qquad \mu_h=\int_\ljacobiR w_h(y)\,dy.
\]
Here $k=4l+1\geq9$ is the original tensor order. These are the
original programme objects, and no existence assertion for an
off-line packet is added. The roots of $h$ are closed under conjugation
with unchanged multiplicities, so $h$ has real coefficients. The
original completed zeta function satisfies
$g(\bar s)=\overline{g(s)}$. Thus $w_h(-y)=w_h(y)$; convolution
preserves this equality and $m_k$ is even. Put
\[
 c=k/2,\quad q=[1+k(m-1)](k+1)^2,\qquad
 \chi(S)=\prod_{a,b=0}^k
 [S-c-(2a-k)\delta-i(2b-k)\gamma]^{1+k(m-1)}.
\]
The full finite quotient will use $N\geq q-1$.

The reference measure and its mass are exactly
\[
 d\sigma(y)=\rho_\sigma(y)\,dy
 =\frac{|\Gamma(1/4+iy/2)|^2}{2\pi}\,dy,
 \qquad M_\sigma=\sqrt{2\pi}.
\]
In particular, this measure is not replaced by a probability measure.
The already proved CAI16--18 envelope supplies
\begin{equation}\label{ljacobi:eq:envelope}
 a_k(1+y^2)^{-M}\rho_\sigma(y)\leq m_k(y)
                       \leq A_k\rho_\sigma(y),\qquad M=21+4m,
\end{equation}
where the retained constants are
\[
 \alpha=\pi/2,\quad p=8m-9/2,\quad B_h=42+8m,\quad
 \vartheta_h=\int_{-1}^1w_h(y)\,dy,
 \quad M_\alpha=\int_\ljacobiR e^{\alpha y}w_h(y)\,dy,
\]
\[
 c_\Gamma=\sqrt2e^{-7/6},\qquad C_\Gamma=\sqrt{2\sqrt5}e^{2/3},
\]
\[
 a_k=\frac{c_h\vartheta_h^{k-3}e^{-\alpha(k-3)}(k-2)^{-B_h}}
              {C_\Gamma2^{B_h/2}},\qquad
 A_k=\frac{C_h^{\rm tilt}}{c_\Gamma}k^{p+1}M_\alpha^{k-1}.
\]
The constants $c_h,C_h^{\rm tilt}$ are the original CAI lower and
upper envelope constants. Their proof is not replaced by a hypothesis:
the source is the retained programme provider
\texttt{CAI\_\allowbreak COMPLETE\_\allowbreak PREREQUISITE.tex},
equations CAI16--18.
The all-degree consequence, with
$\Delta_N=[1+4(N+M)^2]^M$, is
\begin{equation}\label{ljacobi:eq:finite-envelope}
 \frac{a_k}{\Delta_N}\|P\|_\sigma^2
       \leq\int|P(y)|^2m_k(y)\,dy\leq A_k\|P\|_\sigma^2
       \qquad(\deg P\leq N).
\end{equation}
The notation $\Delta_N$ here distinguishes the retained scalar CAI
constant from the source-coordinate matrix $D_N$ below.

The human literature input is Franti\v{s}ek \v{S}tampach and Pavel
\v{S}\v{t}ov\'{i}\v{c}ek, \emph{Spectral representation of some
weighted Hankel matrices and orthogonal polynomials from the Askey
scheme}; see
\href{https://arxiv.org/abs/1811.05820v1}{arXiv:1811.05820v1}.
We use Theorem \texttt{thm:commut\_\allowbreak matr\_\allowbreak Jac},
its following remark and proof,
and Theorem \texttt{thm:def\_H}. The following is a direct reconstruction
for the actual multiplier, including the mass absent from that
source's probability-measure convention.

\section{Gamma polynomials and the native Jacobi operator}

Define the monic polynomials by
\begin{equation}\label{ljacobi:eq:generating}
 G(y,t)=\sum_{n\geq0}p_n(y)\frac{t^n}{n!}
       =(1+t^2)^{-1/4}e^{y\arctan t}.
\end{equation}
Every branch here is its analytic branch at $t=0$. To verify the
mass and norms directly, the beta integral gives, for $a>0$,
\[
 \frac{\Gamma(a+i\tau)\Gamma(a-i\tau)}{\Gamma(2a)}
 =\int_\ljacobiR\frac{e^{i\tau v}}{(2\cosh(v/2))^{2a}}\,dv.
\]
This follows from $u=e^v$ in the beta integral on $(0,\infty)$.
Fourier inversion, followed by analytic continuation in the strip
where the Gamma envelope is integrable, gives at $a=1/4$ and $y=2\tau$
\begin{equation}\label{ljacobi:eq:mgf-sigma}
 \int_\ljacobiR e^{zy}\,d\sigma(y)
 =\sqrt{2\pi}(\cos z)^{-1/2},\qquad |\Re z|<\pi/2.
\end{equation}
In detail, inversion gives the Fourier transform at real frequency
$v$ as $2\Gamma(1/2)/(2\cosh v)^{1/2}$; setting $v=-iz$ yields
$\sqrt{2\pi}(\cos z)^{-1/2}$. The CAI17 Gamma envelope bounds the
integrand on every closed narrower strip and justifies holomorphy
and this continuation. At zero this proves the stated mass.

For sufficiently small $t,u$, the cosine addition formula and
\eqref{ljacobi:eq:mgf-sigma} imply
\[
 \int_\ljacobiR G(y,t)G(y,u)\,d\sigma(y)
 =M_\sigma(1-tu)^{-1/2}
 =M_\sigma\sum_{n\geq0}\frac{(1/2)_n}{n!}t^nu^n.
\]
All coefficient extractions may be taken on fixed small complex
circles, dominated by $Ce^{\varepsilon|y|}$ with $\varepsilon<\pi/2$.
Consequently
\[
 \int p_m p_n\,d\sigma=\delta_{mn}\gamma_n,
 \qquad\gamma_n=M_\sigma n!(1/2)_n.
\]
Define, without altering the measure,
\[
 \varphi_n=\frac{p_n}{\sqrt{\gamma_n}},\qquad
 r_n=\sqrt{M_\sigma}\varphi_n
     =\frac{p_n}{\sqrt{n!(1/2)_n}}.
\]
Thus $\varphi_0=M_\sigma^{-1/2}$, $r_0=1$, and
$\int r_mr_n\,d\sigma=M_\sigma\delta_{mn}$.
The $r_n$ are not orthonormal in $d\sigma$.

Differentiating \eqref{ljacobi:eq:generating} gives
$(1+t^2)\partial_tG=(y-t/2)G$ and hence
\[
 p_{n+1}=yp_n-n(n-1/2)p_{n-1},\qquad p_0=1,\quad p_1=y.
\]
Therefore
\begin{equation}\label{ljacobi:eq:recurrence}
 y\varphi_n=b_n\varphi_{n+1}+b_{n-1}\varphi_{n-1},\qquad
 b_n=\sqrt{(n+1)(n+1/2)},\quad b_{-1}=0,
\end{equation}
and the same recurrence holds for $r_n$. Also $p_n(-y)=(-1)^np_n(y)$.
The original source's Meixner--Pollaczek generating function gives
the exact dictionary
\[
 p_n(y)=n!P_n^{(1/4)}(y/2;\pi/2),\qquad
 r_n(y)=\sqrt{\frac{n!}{(1/2)_n}}P_n^{(1/4)}(y/2;\pi/2).
\]
Both the coordinate $y/2$ and every normalization factor remain.

The polynomials are complete in $L^2(\sigma)$. Indeed, if $f$ is
orthogonal to all polynomials, then $f\,d\sigma$ is a finite complex
measure, and Cauchy--Schwarz together with
\eqref{ljacobi:eq:mgf-sigma} shows that
$\int f(y)e^{zy}\,d\sigma(y)$ is holomorphic for $|\Re z|<\pi/4$.
Every derivative at zero is zero, so analytic uniqueness makes
the transform zero throughout that connected strip. Fourier
uniqueness on the imaginary axis makes $f\,d\sigma=0$, hence $f=0$.
The identical argument with the finite measure $(1+y^2)d\sigma$
proves polynomial density in that weighted space as well.

It follows that
\[
 U_\sigma:\ell^2(\ljacobiN_0)\longrightarrow L^2(\sigma),\qquad
 U_\sigma e_n=\varphi_n
\]
is unitary. The operator $J=U_\sigma^*M_yU_\sigma$ is self-adjoint,
and its matrix $\ljacobiJ$ has the entries \eqref{ljacobi:eq:recurrence}, with zero
diagonal. Polynomial density for $(1+y^2)d\sigma$ means that the
polynomials are a core of $M_y$ in its graph norm. Thus the finite
vectors are a core of $J$, so this is the unique self-adjoint closure
of the stated Jacobi matrix on finite vectors.

\section{The original bounded multiplier and exact first column}

Set $R_k=m_k/\rho_\sigma$. By \eqref{ljacobi:eq:envelope},
\[
 0<a_k(1+y^2)^{-M}\leq R_k(y)\leq A_k.
\]
Therefore multiplication by this exact density ratio is bounded,
positive and self-adjoint on $L^2(\sigma)$, and
\begin{equation}\label{ljacobi:eq:operator}
 \ljacobiCC=U_\sigma^*M_{R_k}U_\sigma,\qquad
 0\leq\ljacobiCC\leq A_k I,\qquad
 \ljacobiCC_{mn}=\int_\ljacobiR\varphi_m(y)\varphi_n(y)m_k(y)\,dy.
\end{equation}
This is not merely a formally defined matrix. The bounded multiplier
preserves $\ljacobiDom M_y$, since
$\|yR_k f\|_2\leq A_k\|yf\|_2$. Consequently $\ljacobiCC$ preserves
$\ljacobiDom J$ and $J\ljacobiCC=\ljacobiCC J$ there. Entrywise the same statement is
\[
 b_{m-1}\ljacobiCC_{m-1,n}+b_m\ljacobiCC_{m+1,n}
 =b_{n-1}\ljacobiCC_{m,n-1}+b_n\ljacobiCC_{m,n+1}.
\]
Each side is a finite sum; it also follows at once by integrating
\eqref{ljacobi:eq:recurrence} against the other polynomial and $m_k$.

Put $\alpha_j=\ljacobiCC_{j0}$. The exact native-mass expression is
\begin{equation}\label{ljacobi:eq:alpha-moment}
 \alpha_j=\frac1{M_\sigma}\int r_j(y)m_k(y)\,dy
 =\frac{\int p_j(y)m_k(y)\,dy}
             {M_\sigma\sqrt{j!(1/2)_j}}.
\end{equation}
In particular $\alpha_0=\mu_h^k/M_\sigma$ by Tonelli.

\begin{ljacobitheorem}\label{ljacobi:thm:first}
For every $m,n\geq0$ the complete entry is the finite sum
\begin{equation}\label{ljacobi:eq:first}
 \ljacobiCC_{mn}=\sum_{j=|m-n|}^{m+n}\alpha_j[r_j(\ljacobiJ)]_{mn}.
\end{equation}
Here $r_j(\ljacobiJ)$ is evaluated in the full Jacobi matrix, not in a
premature finite compression.
\end{ljacobitheorem}
\begin{proof}
The recurrence and $r_0=1$ give $r_n(\ljacobiJ)e_0=e_n$ by induction.
Indeed, for $n=0$ the recurrence gives $r_1(\ljacobiJ)e_0=e_1$; if the
assertion holds at $n,n-1$, then
$b_n r_{n+1}(\ljacobiJ)e_0=\ljacobiJ e_n-b_{n-1}e_{n-1}=b_ne_{n+1}$.
Expand the degree-$m+n$ polynomial $r_mr_n$ in the $r_j$ basis.
Orthogonality shows that its coefficient at $r_j$ is
\[
 d_{mnj}=\frac1{M_\sigma}\int r_mr_nr_j\,d\sigma
       =\int\varphi_m r_j\varphi_n\,d\sigma
       =[r_j(\ljacobiJ)]_{mn}.
\]
The last equality follows either from the multiplication operator
or by applying \eqref{ljacobi:eq:recurrence} finitely many times. Its
integral is zero when $j>m+n$ by orthogonality. It is also zero
when $m+j<n$ or $n+j<m$, by applying orthogonality to $r_n$ or
$r_m$ respectively. Thus only $|m-n|\leq j\leq m+n$ occurs.
Integrate this polynomial identity against $m_k/M_\sigma$ and
use \eqref{ljacobi:eq:alpha-moment}. The result is \eqref{ljacobi:eq:first}.
\end{proof}

Evenness of the unchanged weight gives $\alpha_j=0$ for odd $j$ and
$\ljacobiCC_{mn}=0$ for odd $m+n$. The first column also reconstructs the
actual multiplier in $L^2(\sigma)$, with the exact mass factor:
\begin{equation}\label{ljacobi:eq:mass-caveat}
 U_\sigma\ljacobiCC e_0=\frac{R_k}{\sqrt{M_\sigma}},\qquad
 R_k=\sqrt{M_\sigma}U_\sigma\ljacobiCC e_0
       =\sum_{j\geq0}\alpha_jr_j\quad\text{in }L^2(\sigma).
\end{equation}
To prove the last convergence, expand the vector $\ljacobiCC e_0\in\ell^2$
in its coordinates $\alpha_j$ and apply the unitary $U_\sigma$.
In particular the multiplier itself is not $U_\sigma\ljacobiCC e_0$ in
the original non-probability convention. No pointwise or
operator-norm convergence of the displayed series is asserted.

Define $M_h(z)=\int e^{zy}w_h(y)\,dy$. CAI16 gives
$M_\alpha<\infty$. Since the weight is even,
$\int e^{\alpha|y|}w_h\leq2M_\alpha$. Tonelli and convolution give
\[
 \int e^{zy}m_k(y)\,dy=M_h(z)^k,
 \qquad |\Re z|<\alpha.
\]
For real $z$ this is the absolutely convergent product integral;
the same domination gives the complex version and holomorphy.
For $|t|\leq1/4$,
$|\arctan t|\leq\operatorname{arctanh}(1/4)<\alpha$.
We can therefore integrate \eqref{ljacobi:eq:generating} on any smaller
complex circle and extract coefficients. This proves exactly
\begin{equation}\label{ljacobi:eq:alpha-generating}
 \boxed{\alpha_j=
 \frac{j!}{M_\sigma\sqrt{j!(1/2)_j}}
 [t^j]\left((1+t^2)^{-1/4}M_h(\arctan t)^k\right).}
\end{equation}
The factor $1/M_\sigma=1/\sqrt{2\pi}$ is necessary.

\section{The finite-cutoff sign and its full boundary vector}

Let $P_N$ project to $e_0,\ldots,e_N$, and write
\[
 C_N=P_N\ljacobiCC P_N,\quad J_N=P_N\ljacobiJ P_N,\quad
 c_N=(\ljacobiCC_{0,N+1},\ldots,\ljacobiCC_{N,N+1})^T.
\]
All are now finite coordinate matrices/vectors. Compressing the
full entrywise commutation identity gives
\[
 0=[J_N,C_N]+P_N\ljacobiJ(I-P_N)\ljacobiCC P_N
                    -P_N\ljacobiCC(I-P_N)\ljacobiJ P_N.
\]
The unique crossing Jacobi edge is $b_N$. Since $\ljacobiCC$ is real
symmetric, the two extra terms are respectively $b_Ne_Nc_N^*$
and $b_Nc_Ne_N^*$. Hence
\begin{equation}\label{ljacobi:eq:cutoff}
 \boxed{[J_N,C_N]=b_N(c_Ne_N^*-e_Nc_N^*).}
\end{equation}
This proves the stated sign. Formula \eqref{ljacobi:eq:first} computes
$C_N$ using the coefficients through index $2N$. It computes $c_N$
using those through index $2N+1$. Because the native odd coefficients
vanish, $\alpha_{2N+1}$ is zero; retaining it in the index range is
correct and harmless.

In fact evenness gives $(c_N)_N=\ljacobiCC_{N,N+1}=0$, so $c_N\perp e_N$.
If $c_N=0$, the finite matrices commute. If $c_N\ne0$, the
commutator vanishes on the orthogonal complement of
$\operatorname{span}\{e_N,c_N\}$ and acts on the ordered orthonormal
basis $e_N,c_N/\|c_N\|$ there as
\[
 \begin{pmatrix}0&-b_N\|c_N\|\\b_N\|c_N\|&0\end{pmatrix}.
\]
Thus its rank is two and its norm is exactly $b_N\|c_N\|$.
These are identities, not new estimates on the vector $c_N$.

The distinction between $P_Nr_j(\ljacobiJ)P_N$ and $r_j(J_N)$ is essential.
For example $r_2(y)=(y^2-1/2)/\sqrt{3/2}$, and at $N=1$,
\[
 [r_2(\ljacobiJ)]_{11}=\sqrt6,\qquad [r_2(J_1)]_{11}=0.
\]
Indeed $(\ljacobiJ^2)_{11}=b_0^2+b_1^2=1/2+3$, whereas
$(J_1^2)_{11}=b_0^2=1/2$. The missing term is the path from $1$
to $2$ and back to $1$. Accordingly
$\ljacobiCC_{11}=\alpha_0+\sqrt6\alpha_2$, not merely $\alpha_0$.

\section{Exact original source and observation coordinates}

For $F(S)=\sum_{j=0}^N a_jS^j$, the monomial coefficient vector
of the same function $F(c+iy)$ is $Ta$, where
\[
 T_{rj}=\binom jr c^{j-r}i^r\ (r\leq j),\qquad T_{rj}=0\ (r>j).
\]
This is the binomial theorem with no conjugation or mass change.
Let $B$ have the monomial coefficient vector of $\varphi_n$ as
column $n$. It is upper triangular with diagonal
$1/\sqrt{\gamma_n}$, while $T$ has diagonal $i^n$.
Both matrices are invertible. If $x$ denotes the Gamma-frame
coefficient vector of $F(c+iy)$, equality of the two polynomial
expansions is $Bx=Ta$. Hence
\begin{equation}\label{ljacobi:eq:source}
 x=D_Na,\qquad D_N=B^{-1}T,\qquad
 (D_N)_{nn}=i^n\sqrt{\gamma_n}.
\end{equation}
Integration in the unchanged arithmetic measure gives
\begin{equation}\label{ljacobi:eq:gram}
 H_{S,N}^{\rm ar}=D_N^*C_ND_N.
\end{equation}
For every $x\in\ljacobiC^{N+1}$, \eqref{ljacobi:eq:finite-envelope} applied to
$\sum x_n\varphi_n$ gives the concrete positive bounds
\begin{equation}\label{ljacobi:eq:CN-bounds}
 \frac{a_k}{\Delta_N}I\leq C_N\leq A_kI.
\end{equation}
Thus the inverse in the next formula exists for every fixed $N$.
No positive lower bound independent of $N$ for the infinite
operator is inferred from the decaying lower density envelope.

Retain the actual observation map
$\mathcal A_S=\Lambda\mathcal J_{S,N}$, including its full physical
Taylor unit, in its original source and observation coordinates.
Its codomain is the original observation image with its retained
Hermitian structure, so it is onto. Its matrix in the Gamma-frame
source coordinates must be
\begin{equation}\label{ljacobi:eq:E}
 \mathcal E_N=\mathcal A_SD_N^{-1},
\end{equation}
since $\mathcal A_Sa=\mathcal A_SD_N^{-1}x$. The symbol
$\mathcal J_{S,N}$ denotes the original class map, not the Jacobi
matrix $J_N$.

\begin{ljacobitheorem}\label{ljacobi:thm:quotient}
For every observed vector $u$ the original attained minimum is
\begin{equation}\label{ljacobi:eq:quotient}
 \min_{\mathcal A_Sa=u}a^*H_{S,N}^{\rm ar}a
  =u^*Q_{B,N}u,\qquad
 Q_{B,N}=(\mathcal E_NC_N^{-1}\mathcal E_N^*)^{-1},
\end{equation}
and its unique minimizing source vector is
\begin{equation}\label{ljacobi:eq:lift}
 a_{\min}=D_N^{-1}C_N^{-1}\mathcal E_N^*Q_{B,N}u.
\end{equation}
\end{ljacobitheorem}
\begin{proof}
Surjectivity of $\mathcal E_N$ makes its adjoint injective, so
$K=\mathcal E_NC_N^{-1}\mathcal E_N^*$ is strictly positive on the
observation space. Put $x_0=C_N^{-1}\mathcal E_N^*K^{-1}u$.
Then $\mathcal E_Nx_0=u$ and
$x_0^*C_Nx_0=u^*K^{-1}u$. Every other lift is $x_0+z$ with
$\mathcal E_Nz=0$, and
$z^*C_Nx_0=(\mathcal E_Nz)^*K^{-1}u=0$. Therefore its energy is
$u^*K^{-1}u+z^*C_Nz$, strictly larger unless $z=0$.
The bijection $x=D_Na$ and \eqref{ljacobi:eq:gram} prove both formulas in
the original source coordinates. Equivalently,
\[
 H_{S,N}^{{\rm ar},-1}=D_N^{-1}C_N^{-1}(D_N^*)^{-1},\qquad
 \mathcal A_SH_{S,N}^{{\rm ar},-1}\mathcal A_S^*
                   =\mathcal E_NC_N^{-1}\mathcal E_N^*.
\]
Thus this is exactly the original quotient metric, not an auxiliary
minimum for a changed source.
\end{proof}

For completeness, \eqref{ljacobi:eq:CN-bounds} also gives the exact matrix
comparison
\[
 \frac{a_k}{\Delta_N}(\mathcal E_N\mathcal E_N^*)^{-1}
       \leq Q_{B,N}\leq
 A_k(\mathcal E_N\mathcal E_N^*)^{-1}.
\]
This follows by inverting the two positive bounds for $C_N$,
congruence with $\mathcal E_N$, and inversion once again. It does
not replace $\mathcal E_N$ by another observation map.

\section{Explicit finite checks and retained phase problem}

Write $\nu_j=\int y^jm_k(y)\,dy$ and
$\eta_j=\int y^jw_h(y)\,dy$, so $\eta_0=\mu_h$.
The first polynomial values from the exact recurrence are
\[
 p_0=1,\quad p_1=y,\quad p_2=y^2-\tfrac12,\quad
 p_3=y^3-\tfrac72y,\quad p_4=y^4-11y^2+\tfrac{15}4.
\]
Consequently the first even column entries are
\[
 \alpha_0=\frac{\nu_0}{M_\sigma},\qquad
 \alpha_2=\frac{\nu_2-\nu_0/2}{M_\sigma\sqrt{3/2}},\qquad
 \alpha_4=\frac{\nu_4-11\nu_2+15\nu_0/4}
                   {M_\sigma\sqrt{315/2}}.
\]
Multinomial expansion of the actual convolution, using evenness,
gives
\[
 \nu_0=\mu_h^k,\quad \nu_2=k\eta_2\mu_h^{k-1},\quad
 \nu_4=k\eta_4\mu_h^{k-1}
              +3k(k-1)\eta_2^2\mu_h^{k-2}.
\]
The $3k(k-1)$ factor counts the two distinct factors contributing
second powers: $\binom{k}{2}\binom42=3k(k-1)$.

At cutoff $N=1$, evenness gives
\[
 C_1=\frac1{M_\sigma}\begin{pmatrix}\nu_0&0\\0&2\nu_2\end{pmatrix},
 \quad J_1=\begin{pmatrix}0&1/\sqrt2\\1/\sqrt2&0\end{pmatrix},
 \quad c_1=\begin{pmatrix}\alpha_2\\0\end{pmatrix},\quad b_1=\sqrt3.
\]
Its $(0,1)$ commutator entry is
$(2\nu_2-\nu_0)/(\sqrt2M_\sigma)=\sqrt3\alpha_2$,
confirming both the sign and the mass factor in \eqref{ljacobi:eq:cutoff}.
The degree-two source change is explicitly
\[
 D_2=\sqrt{M_\sigma}
 \begin{pmatrix}
 1&c&c^2-1/2\\
 0&i/\sqrt2&\sqrt2ic\\
 0&0&-\sqrt{3/2}
 \end{pmatrix}.
\]
In particular the source constant $F=1$ has coefficient
$x_0=\sqrt{M_\sigma}$ and energy
$M_\sigma(C_0)_{00}=\mu_h^k$, retaining the original arithmetic mass.

The companion exact symbolic script \texttt{check\_symbolic.py}
checks \eqref{ljacobi:eq:first} for all $0\leq m,n\leq4$ (25 entries),
checks \eqref{ljacobi:eq:cutoff} at $N=0,1,2,3$, checks the generating
polynomials through degree six, and checks
\eqref{ljacobi:eq:source}--\eqref{ljacobi:eq:gram} through degree two. All moment
variables are independent real symbols; the checks even allow
nonzero odd moments. The parameter representing $M_\sigma$ remains
a symbolic positive mass rather than being set to one. These
finite checks supplement, and do not replace, the general proofs.

For the original classes $a_1=[p_N^{\rm ar}]_\chi$,
$a_2=[p_{N+1}^{\rm ar}]_\chi$ and original terminal class $v$, the
retained quantities are therefore exactly
\[
 B_i=(\Lambda a_i)^*Q_{B,N}\Lambda v,\qquad
 \Phi_{B,N}=\frac{2\Re(\overline{B_1}B_2)}{\omega_{N,k}^{\rm ar}}.
\]
The formulas above give a finite representation of the same
quantities and the full Jacobi compression defect. They do not
evaluate or estimate these complex signed products. The physical
Taylor unit, the observation map, and both original classes remain
present; parity of $C_N$ does not remove them.


\endgroup

% END INDEXED LITERATURE EXACT RECEIVERS 20260920

\end{document}
